Short-horizon architecture and optimizer comparisons are the default currency of empirical
language-model research. We report a controlled attempt to replicate one such comparison ---
our own --- and show that the headline result was a property of the measurement recipe rather
than of the architectures being compared.

A single-seed run on an RTX 3070 Ti at batch size 8 placed the point where softmax attention
overtakes minGRU between \textbf{6.6M and 7.4M} training tokens on a 12-layer / 768-dim model.
Re-running the same matched pair on an NVIDIA GH200 at \textbf{n = 5 seeds} through a 50M-token
horizon does not reproduce that location. Instead the mean validation-loss curves cross
\textbf{twice}: minGRU overtakes attention at \textbf{1.05M} tokens and attention overtakes minGRU for
good at \textbf{12.35M} tokens, holding a \textbf{-0.227} {[}-0.244, -0.211{]} advantage at 50M. Within the
original 6.6--8.2M window minGRU leads by \textbf{0.13--0.18} nats on every seed.

We then isolate the confounds one at a time, each with five seeds. Truncating the run to 20M
tokens while letting the cosine schedule follow the truncated step count moves the late crossing
from 12.34M to \textbf{14.58M}; restoring the 50M cosine horizon and stopping at the same 20M budget
recovers \textbf{12.34M} (per-seed 11.93--12.58M) and reproduces the 19.7M gap to within 0.001. Holding
batch size at the original 8 on the new hardware produces \textbf{no crossing at all} through 7.38M
tokens on any seed. A matched-batch 50M ranking of ten mixer configurations places attention
first at \textbf{4.222} {[}4.204, 4.240{]} but statistically tied with a 10$\times$minGRU + 2$\times$attention hybrid at
\textbf{4.232} {[}4.210, 4.254{]}.

We observe the same phenomenon on an independent axis. In a five-stage optimizer funnel on Apple
silicon, MONA ranks first at 16M parameters and \textbf{last of four} at the \texttt{exact\_128m} configuration (128,367,988 parameters) --- ``exact'' there names the funnel's bit-exactness resume gate, recorded in
\path{research/exact-128m-gate-polar.json}, not a round parameter count;
the eventual 128M champion, Polar Express Muon, ranked third in the initial learning-rate screen
and loses to Muon NS5 when run back down at small scale. We also correct the confidence intervals
originally computed for that funnel: with the correct two-seed Student-\emph{t} multiplier, the
champion-selection intervals \textbf{overlap}, and the selection rests on a sign-consistent 2-of-2
ordering rather than on separated intervals.

The transferable claim is not ``attention wins late.'' It is that \textbf{the crossing token is a
function of batch size, learning-rate horizon, hardware and evaluation protocol}, and that a
ranking measured at one recipe carries no license to be reported as a property of the
architectures.

One limitation is large enough to state in the abstract. Every run here uses standard
parametrization with a single global Muon learning rate, so we cannot yet separate ``the ranking
depends on the recipe'' from ``the arms were unequally well tuned at each recipe.'' \S8 states that
competing explanation, gives the evidence that partially bounds it, and specifies the $\mu$P arm ---
a 2 $\times$ 2 over parametrization and architecture on the suite 24 recipe, plus two cheaper ablations ---
that would settle it, with a pre-registered reading of each possible outcome.

\begin{center}\rule{0.5\linewidth}{0.5pt}\end{center}

\section{1. Introduction}\label{introduction}

A large fraction of practical architecture and optimizer selection happens on short runs. The
reasoning is economic: a screen that costs 1\% of a full run and preserves the ordering of
candidates is worth more than the full run. The assumption doing the work is that the ordering is
preserved.

That assumption has been challenged before, along several independent axes, and it has usually
lost. Melis et al.~re-evaluated a sequence of LSTM successors under large-scale black-box tuning
and found that properly regularized standard LSTMs outperformed the newer architectures that had
been reported to beat them \cite{ref6}. Musgrave et al.~audited four years of deep metric learning and
found the accumulated gains to be marginal once experimental methodology was controlled \cite{ref38}.
Henderson et al.~showed that non-determinism and intrinsic variance in deep reinforcement learning
make published comparisons difficult to interpret without seed discipline and significance
reporting \cite{ref8}. Dodge et al.~made the dependence explicit: test-set scores alone are insufficient
for model comparison, and several published comparisons \textbf{invert} when the computation budget
spent on hyperparameter search changes \cite{ref7}. Hooker's \emph{hardware lottery} names the systems version
of the same problem --- an idea can win because it suits the available hardware and software rather
than because it is better \cite{ref26}.

The specific variable this paper is about --- \textbf{training horizon} --- is less often treated as a
ranking confound, even though the pieces needed to predict that it would be are individually well
established. Optimal learning rate depends strongly on token horizon, with longer training
requiring smaller learning rates and the optimal value following its own scaling law \cite{ref31}. The
number of steps required to reach a target loss varies enormously with batch size and workload,
and much of the historical disagreement about whether batch size affects final quality is
attributable to differences in metaparameter tuning at different batch sizes \cite{ref27}. The largest
useful batch size is itself a moving target predicted by the gradient noise scale, which changes
\emph{over the course of a run} as loss decreases \cite{ref28}, and critical batch size scales primarily with
data size rather than model size \cite{ref29}. Small batches are not merely a scaled-down version of large
ones: they require different optimizer hyperparameters, and holding Adam's second-moment decay
fixed rather than its token half-life mis-specifies them \cite{ref30}. Compute-optimal scaling work
established that conclusions drawn at fixed data are systematically misleading \cite{ref9,ref10}, and
Choshen et al.~found enough seed-level variability in scaling-law estimation that training several
small models is sometimes more informative than training one large one \cite{ref37}.

Put together, these say that batch size, learning-rate schedule and horizon jointly determine
where a training curve is at any given token count. What they do not say --- and what we could not
find measured directly --- is what that does to the \textbf{ordering} of two architectures whose curves
cross, and specifically to the \emph{token at which they cross}, which is the quantity that short-run
comparisons implicitly report.

This paper reports what happened when we tested that against a result we had previously published
ourselves.

\textbf{The original observation.} In June 2026 we ran three 12-layer / 768-dim mixers --- minGRU,
softmax attention, and Mamba-2 \cite{ref1,ref2} --- on FineWeb-Edu \cite{ref25} for 8.192M tokens on an RTX 3070 Ti
Laptop (batch size 8, sequence length 512, Muon, one seed). minGRU led at every early checkpoint;
attention closed the gap monotonically and overtook between 6.6M and 7.4M tokens (suite 14). Two
follow-up pairs at a smaller model and a longer horizon preserved the \emph{shape} --- early recurrent
lead, late attention overtake --- with the flip landing near 8.19M and 7.37M tokens respectively
(suite 15). We described this as ``bias wins early, capacity wins late'' and treated the \textasciitilde7M token
location as the finding.

That interpretation was not unreasonable in isolation. It is consistent with the broader picture
in which pure state-space and recurrent stacks match Transformers on many tasks but lag on those
requiring copying or in-context retrieval \cite{ref5}, and with the general expectation that an
architecture's inductive bias should matter most when data is scarce. What made it a mistake was
not the story. It was that we reported a \textbf{number} --- the crossing token --- from a single seed on a
single recipe, and treated it as a property of the architectures.

\textbf{What this paper adds.} Access to an NVIDIA GH200 made a five-seed replication affordable. We
ran the matched attention/minGRU pair plus eight additional mixer configurations to 50M tokens,
then ran three targeted isolates that each change exactly one factor. The specific contributions:

\begin{enumerate}
\def\labelenumi{\arabic{enumi}.}
\tightlist
\item
  \textbf{A failed location replication with a measured alternative.} At n = 5 on the GH200 the
  6.6--7.4M crossing does not appear; two crossings appear instead, at 1.05M and 12.35M
  (\S4.2). Every seed agrees on the sign at every marker except the 12.3M tie.
\item
  \textbf{Learning-rate horizon is a first-order confound for the crossing token, not only for final
  loss.} Truncating the token budget without holding the cosine horizon fixed moves the late
  crossing by 2.248M tokens --- 18\% of its location --- and the effect is measurable in the learning
  rate itself, which at the crossing marker is \textasciitilde47\% of its long-horizon value (\S4.3). This is the
  ranking-level consequence of the horizon-dependence of optimal LR reported in \cite{ref31}.
\item
  \textbf{Batch size, not just hardware, changes the early ordering.} At batch size 32 attention
  leads the first evaluation on all five seeds; at batch size 8 on the \emph{same} GPU minGRU leads
  the first evaluation on all five seeds and never surrenders it through 7.38M tokens (\S4.4).
  Both batch sizes are far below any plausible critical batch size at this scale \cite{ref29}, so this is
  not a large-batch degradation effect.
\item
  \textbf{A matched-batch 50M mixer ranking} in which the previously reported ordering of hybrids is
  substantially rearranged once batch and evaluation count are held fixed (\S4.5), and in which
  the best hybrid is statistically indistinguishable from pure attention.
\item
  \textbf{An independent replication of the phenomenon on the optimizer axis}, including a
  statistical correction to our own previously reported selection intervals (\S5). Zhao et al.
  found that most modern optimizers perform comparably at their own optimal hyperparameters
  \cite{ref36}; our funnel is a case study in what happens when the tuning budget is too short to find
  those hyperparameters.
\item
  \textbf{A protocol} (\S9) and an explicit list of claims we withdraw (\S7.3).
\item
  \textbf{A named, designed follow-up} (\S8): the single most likely explanation we did \emph{not} control
  for is parametrization, and we specify the $\mu$P arm that would test it.
\end{enumerate}

We deliberately do not claim a mechanism. ``Recurrent inductive bias helps early, attention
capacity wins late'' remains a plausible story that fits these curves; it is not established by
them, and it does not predict the crossing token, which is the quantity that actually moved.

\begin{center}\rule{0.5\linewidth}{0.5pt}\end{center}

\section{2. Background and related work}\label{background-and-related-work}

\textbf{Sequence mixers.} minGRU is a minimal gated recurrent unit whose recurrence is
parallelizable over sequence length, proposed as evidence that heavily simplified RNNs remain
competitive with Transformers at moderate scale \cite{ref1}. Mamba-2 recasts selective state-space models
and attention variants inside a shared structured-semiseparable-matrix framework, yielding a
core layer 2--8$\times$ faster than Mamba while remaining competitive on language modeling \cite{ref2}. Gated
DeltaNet combines gated adaptive memory control with the delta update rule to improve
retrieval-limited behavior in linear Transformers \cite{ref3}. Multi-head Latent Attention compresses the
KV cache into a latent vector and was introduced with DeepSeek-V2 \cite{ref4}. All four, plus attention
and six hybrid interleavings, are implemented in the \texttt{nanolab} codebase used here.

\textbf{Hybrids.} Interleaving a small number of attention layers into an otherwise recurrent or
state-space stack is now a standard recipe. At 8B parameters and 3.5T tokens, a 43\% Mamba-2 / 7\%
attention / 50\% MLP hybrid exceeded a matched Transformer on all twelve standard tasks evaluated,
while pure SSMs lagged on copying and in-context-learning tasks \cite{ref5}. Our \S4.5 board is a
small-scale, matched-recipe version of the same question, and reaches a compatible conclusion:
the best hybrid is statistically indistinguishable from pure attention, and pure recurrent stacks
are clearly behind.

\textbf{Benchmark validity and budget-dependent conclusions.} The methodological literature
summarized in \S1 \cite{ref6,ref7,ref8,ref38} establishes that measured rankings depend on the resources spent
measuring them. That work is almost always framed around \emph{hyperparameter search} budget or
implementation variance. The closest precedent to our framing is Dodge et al.'s expected
validation performance as a function of budget, which makes the budget an explicit axis of the
reported result rather than an implicit constant \cite{ref7}. We adopt the same posture for the training
horizon: if two curves cross, the crossing token is part of the result and must be reported with
the recipe that produced it.

\textbf{Horizon, batch size and schedule.} Optimal learning rate depends on the token horizon, with
longer runs requiring smaller learning rates and the optimal value obeying its own scaling law
that permits transfer from short horizons to long \cite{ref31}. This is the mechanism behind \S4.3: when a
run is truncated and the cosine schedule follows the truncated step count, the two runs are on
materially different learning rates at every shared token marker, and we measure the ratio at
\textasciitilde47\% at the crossing point. On the batch axis, Shallue et al.~characterized steps-to-target as a
function of batch size across 35 workloads and attributed much of the literature's disagreement
about batch-size effects on quality to differences in metaparameter tuning \cite{ref27}; McCandlish et al.
introduced the gradient noise scale as a predictor of the largest useful batch size and noted that
it \emph{increases as loss decreases within a run} \cite{ref28}; Zhang et al.~measured how critical batch size
scales in pre-training from 85M to 1.2B parameters and found it scales primarily with data size
rather than model size \cite{ref29}. Marek et al.~revisited very small batches and proposed holding Adam's
second-moment half-life fixed in tokens rather than holding its decay rate fixed, finding small
batches stable and more robust to hyperparameter choice than conventional wisdom holds \cite{ref30}. Our
batch-8 and batch-32 recipes (\S4.4) sit well below any critical batch size at this scale, which is
why we read the batch effect there as an ordering effect rather than a large-batch quality
penalty.

\textbf{Parametrization and hyperparameter transfer.} $\mu$P makes many optimal hyperparameters
approximately invariant to width, enabling a tune-small / transfer-large workflow \cite{ref12}; Depth-$\mu$P
extends the analysis to depthwise transfer and identifies fundamental limitations when each
residual block is itself deep, as in modern Transformers \cite{ref33}; u-$\mu$P combines $\mu$P with unit scaling
so that default values are near-optimal and FP8 training works out of the box \cite{ref34}. The picture is
not settled in $\mu$P's favor: Everett et al.~trained tens of thousands of models across four
parameterizations, three optimizers and fourteen model sizes to 26.8B parameters, and found that
\emph{all} parameterizations can achieve hyperparameter transfer, with a per-layer learning-rate
prescription for standard parametrization outperforming $\mu$P \cite{ref32}. Kalra and Barkeshli localized
$\mu$P's apparent advantage over standard parametrization under AdamW to a single factor --- maximizing
the embedding-layer learning rate, which in standard parametrization acts as a bottleneck that
induces instability \cite{ref35}. Our runs use standard parametrization with a global Muon learning rate
and no per-layer prescription; \S8 is about what that leaves untested.

\textbf{Scale as a confound.} Compute-optimal scaling established that model size and token count must
be scaled together and that conclusions at fixed data mislead \cite{ref9,ref10}. Wortsman et al.~reproduced
large-scale instabilities in small models at high learning rates and studied how warm-up, weight
decay and $\mu$Param change loss sensitivity to learning rate \cite{ref11} --- the same class of interaction we
measure at the ranking level in \S4.3. Choshen et al.~derived best practices for estimating scaling
laws from 485 pretrained models and over 1000 fitted laws, finding that fitting to intermediate
checkpoints substantially improves accuracy and that seed variability can make several small
models more useful than one large one \cite{ref37}. Schaeffer et al.~argued that apparently sharp
scale-dependent transitions can be artifacts of the chosen metric \cite{ref13}; our crossings are measured
on a continuous metric (held-out cross-entropy), which removes that particular explanation but not
the recipe dependence.

\textbf{Optimizers.} The optimizer arm of this work compares Muon-family matrix-orthogonalizing
methods against full-parameter baselines. Muon has been shown to scale to large LLM training with
weight decay and per-parameter update-scale adjustment, reporting roughly 2$\times$ computational
efficiency versus AdamW under compute-optimal training \cite{ref14}. Polar Express replaces Newton--Schulz
with a minimax-optimal polynomial iteration for the polar decomposition and reports consistent
validation-loss improvements when integrated into Muon on GPT-2-scale FineWeb training \cite{ref15}.
NorMuon adds neuron-wise adaptive second moments on top of orthogonalized updates \cite{ref16}. MONA adds
a curvature-aware Nesterov acceleration term inside Muon's gradient pipeline and reports gains at
1B--68B MoE scale \cite{ref17}. Muown treats the row-magnitude vector as an explicit optimizer variable and
reports improved perplexity over Muon, SOAP, AdamW and Lion on FineWeb-Edu from 124M to 2.7B \cite{ref18}.
Baselines and controls include AdamW \cite{ref19}, Lion \cite{ref20}, Sophia \cite{ref21}, Schedule-Free AdamW \cite{ref22}, and
Prodigy \cite{ref23}. Two results frame \S5. Zhao et al.~compared SGD, Adafactor, Adam, Lion and Sophia
across model sizes and hyperparameter settings and found that --- SGD aside --- they perform
comparably both at their optima and in robustness to hyperparameter misspecification, suggesting
optimizer choice can be driven by memory and implementation considerations \cite{ref36}. Community
benchmarking efforts such as AlgoPerf exist precisely because optimizer comparisons are so
sensitive to tuning protocol \cite{ref24}. If optimizers are largely comparable \emph{when each is properly
tuned}, then a short screen that cannot tune them is measuring something else --- which is what \S5.4
reports.

\textbf{Data.} All runs in this paper train on FineWeb-Edu, the educational subset of FineWeb \cite{ref25}.

\begin{center}\rule{0.5\linewidth}{0.5pt}\end{center}

\section{3. Experimental setup}\label{experimental-setup}

\subsection{3.1 Model and training}\label{model-and-training}

The mixer experiments use a fixed 12-layer / 768-dim decoder (``124M shape'') with sequence length
512, trained with Muon at learning rate 6e-4 on FineWeb-Edu, \texttt{mixer\_chunk\ =\ 32}, \texttt{compile\ =\ False},
GPU-resident data, and \texttt{eval\_train\ =\ False}. The only lever that varies across arms is the layer
mixer. Seeds are \texttt{1337,\ 42,\ 100,\ 2026,\ 777} for every five-seed suite.

Hardware for suites 22--26 is a Lambda Cloud GH200 (97,871 MiB HBM, aarch64, PyTorch 2.7.0 +
CUDA 12.8). Suites 14--16 ran earlier on an RTX 3070 Ti Laptop (8 GB). \texttt{torch.compile} stalled in
Inductor for attention on GH200 aarch64, so every job in suites 22--26 ran uncompiled; this is a
uniform choice across arms, not a per-arm difference.

Warm-up, evaluation and checkpoint cadences are expressed in \textbf{tokens}, not steps, so that a
wider microbatch still lands evaluations near the original suite-14 markers (batch 8 $\times$ sequence
512 = 4,096 tokens/step; evaluation every 200 such steps, i.e.~every 0.8192M tokens).

\subsection{3.2 Evaluation and statistics}\label{evaluation-and-statistics}

Validation loss is reported at the nearest logged evaluation to each named token marker, with
\texttt{eval\_iters\ =\ 20} for every arm in the matched suites. For a marker with five seeds we report the
mean and a 95\% Student-\emph{t} interval with four degrees of freedom (t$_4$ = 2.776). The gap is defined
as attention - minGRU, so a negative gap means attention is better.

Crossing points are obtained by linear interpolation of the two \textbf{mean} curves on attention's
evaluation grid; per-seed crossings are computed the same way on individual curves and reported as
a range so that a mean crossing cannot be mistaken for a single-seed artifact.

Three rules were applied throughout and are worth stating because they cost us results:

\begin{itemize}
\tightlist
\item
  \textbf{No cross-recipe tables.} Cells measured at batch 8 never appear in the same table as cells
  measured at batch 32, even when the token markers match.
\item
  \textbf{A \texttt{best\_val} field is not a paired snapshot.} Minimum-over-all-evaluations is not a ranking
  and is never reported as one.
\item
  \textbf{An auto-generated intersection over mixed token grids is not a source of truth.} When arms
  ran at different batch sizes their eval grids do not align, and a naive intersection produces
  flat, wrong early markers.
\end{itemize}

\subsection{3.3 Suites}\label{suites}

{%
{\small\begin{longtable}[]{@{}
  >{\raggedright\arraybackslash}p{(\linewidth - 8\tabcolsep) * \real{0.2000}}
  >{\raggedright\arraybackslash}p{(\linewidth - 8\tabcolsep) * \real{0.2000}}
  >{\raggedright\arraybackslash}p{(\linewidth - 8\tabcolsep) * \real{0.2000}}
  >{\raggedright\arraybackslash}p{(\linewidth - 8\tabcolsep) * \real{0.2000}}
  >{\raggedright\arraybackslash}p{(\linewidth - 8\tabcolsep) * \real{0.2000}}@{}}
\toprule\noalign{}
\begin{minipage}[b]{\linewidth}\raggedright
Suite
\end{minipage} & \begin{minipage}[b]{\linewidth}\raggedright
Question
\end{minipage} & \begin{minipage}[b]{\linewidth}\raggedright
Hardware / batch
\end{minipage} & \begin{minipage}[b]{\linewidth}\raggedright
Seeds
\end{minipage} & \begin{minipage}[b]{\linewidth}\raggedright
Budget
\end{minipage} \\
\midrule\noalign{}
\endhead
\bottomrule\noalign{}
\endlastfoot
14 & Does the 2M-token mixer ranking survive to 8.2M? & 3070 Ti / bs8 & 1 & 8.192M \\
15 & Does the crossover shape survive a smaller model and a longer run? & 3070 Ti / bs16 & 1 & 16.4M, 30.3M \\
22 & Does the 6.6--7.4M location replicate at n = 5 to 50M? & GH200 / bs32 (+bs96 drift) & 5 & 50M \\
23 & Independent 20M pair, cosine following \texttt{max\_steps} & GH200 / bs32 & 5 & 20M \\
24 & Same, with the 50M cosine horizon held & GH200 / bs32 & 5 & 20M \\
25 & Original batch size 8, on the new hardware & GH200 / bs8 & 5 & 8.192M \\
26 & Matched-batch 50M ranking of the eight drifted arms & GH200 / bs32 & 5 & 50M \\
\end{longtable}}
}

Suites 24, 25 and 26 are \emph{isolates}: each changes exactly one factor relative to suite 22 and
verifies the recipe fingerprint (\texttt{batch\_size}, \texttt{eval\_iters}, \texttt{max\_steps}, \texttt{lr\_max\_steps}) on every
individual job config before the results are read.

\begin{center}\rule{0.5\linewidth}{0.5pt}\end{center}

\section{4. Results: the mixer axis}\label{results-the-mixer-axis}

\subsection{4.1 The original observation (suite 14)}\label{the-original-observation-suite-14}

RTX 3070 Ti, batch 8, sequence 512, 2,000 steps = 8.192M tokens, seed 1337, one seed per arm.

{%
{\small\begin{longtable}[]{@{}lllll@{}}
\toprule\noalign{}
Tokens & minGRU & Attention & Mamba-2 & gap (A-G) \\
\midrule\noalign{}
\endhead
\bottomrule\noalign{}
\endlastfoot
0.8M & \textbf{6.334} & 6.516 & 6.493 & +0.182 \\
4.1M & \textbf{5.549} & 5.624 & 5.769 & +0.075 \\
6.6M & \textbf{5.353} & 5.358 & 5.560 & +0.005 (tied) \\
7.4M & 5.249 & \textbf{5.239} & 5.469 & -0.010 $\leftarrow$ overtakes \\
8.2M & 5.155 & \textbf{5.136} & 5.383 & -0.019 \\
\end{longtable}}
}

Suite 15 preserved the shape at two other configurations: a 6L/384d model at batch 16 crossed near
\textbf{8.19M} tokens (step 1000 of 2000), and the full 12L/768d model at batch 16 crossed near
\textbf{7.37M} tokens (step 900 of 3700). Each of those four arms is a single seed.

The crossing locations above are read from paired evaluation curves. The end-of-run figures for
these four arms --- 4.852 / 5.021 at 16.4M and 4.260 / 4.575 at 30.3M --- are \texttt{best\_val} fields, i.e.
minimum-over-all-evaluations, and by the rule we set out in \S3.2 they are \textbf{not} paired snapshots
and \textbf{not} a ranking. An earlier version of this section reported them as though they were, which
is exactly the failure \S3.2 is meant to prevent. We state them here only as run summaries, and no
ordering in this paper depends on them. The same caveat applies to the 8.2M row of the suite-14
table above, which is that run's \texttt{done} record rather than a matched-token evaluation.

This is the entire evidential basis for the ``\textasciitilde7M crossover'' claim: \textbf{one seed} at the headline
configuration, and two single-seed pairs whose crossings land at 8.19M and 7.37M --- already a
1.1M-token spread that we under-weighted at the time.

\subsection{4.2 The n = 5 replication (suite 22)}\label{the-n-5-replication-suite-22}

Attention and minGRU share one recipe from step 0: batch 32, \texttt{eval\_iters\ =\ 20}, 50M-token cosine,
five seeds. Means are nearest-evaluation; intervals are 95\% Student-\emph{t} on five seeds.

{%
{\small\begin{longtable}[]{@{}
  >{\raggedright\arraybackslash}p{(\linewidth - 6\tabcolsep) * \real{0.2500}}
  >{\raggedright\arraybackslash}p{(\linewidth - 6\tabcolsep) * \real{0.2500}}
  >{\raggedright\arraybackslash}p{(\linewidth - 6\tabcolsep) * \real{0.2500}}
  >{\raggedright\arraybackslash}p{(\linewidth - 6\tabcolsep) * \real{0.2500}}@{}}
\toprule\noalign{}
\begin{minipage}[b]{\linewidth}\raggedright
Actual tokens
\end{minipage} & \begin{minipage}[b]{\linewidth}\raggedright
Attention {[}95\%{]}
\end{minipage} & \begin{minipage}[b]{\linewidth}\raggedright
minGRU {[}95\%{]}
\end{minipage} & \begin{minipage}[b]{\linewidth}\raggedright
gap A-G {[}95\%{]}
\end{minipage} \\
\midrule\noalign{}
\endhead
\bottomrule\noalign{}
\endlastfoot
0.836M & 6.775 {[}6.755, 6.794{]} & 6.831 {[}6.819, 6.844{]} & \textbf{-0.056} {[}-0.065, -0.048{]} \\
4.112M & 5.888 {[}5.863, 5.914{]} & 5.661 {[}5.637, 5.684{]} & \textbf{+0.228} {[}+0.216, +0.239{]} \\
6.570M & 5.581 {[}5.573, 5.589{]} & 5.402 {[}5.390, 5.414{]} & \textbf{+0.179} {[}+0.169, +0.189{]} \\
7.389M & 5.504 {[}5.489, 5.519{]} & 5.350 {[}5.330, 5.370{]} & \textbf{+0.154} {[}+0.145, +0.162{]} \\
8.208M & 5.434 {[}5.405, 5.462{]} & 5.299 {[}5.267, 5.332{]} & \textbf{+0.134} {[}+0.126, +0.143{]} \\
12.304M & 5.122 {[}5.104, 5.141{]} & 5.120 {[}5.097, 5.144{]} & +0.002 {[}-0.009, +0.013{]} (tie) \\
19.677M & \textbf{4.743} {[}4.733, 4.754{]} & 4.936 {[}4.924, 4.947{]} & \textbf{-0.193} {[}-0.201, -0.185{]} \\
49.988M & \textbf{4.222} {[}4.204, 4.240{]} & 4.449 {[}4.423, 4.476{]} & \textbf{-0.227} {[}-0.244, -0.211{]} \\
\end{longtable}}
}

Interpolated mean-curve crossings: \textbf{1.049M} (minGRU overtakes attention) and \textbf{12.353M}
(attention overtakes minGRU and stays ahead through 50M). Per-seed late crossings are
\textbf{12.031, 12.216, 12.453, 12.474, 12.579M} --- a 0.55M band, so the mean crossing is not one lucky
seed. The sign of the gap agrees across all five seeds at every marker except the 12.3M tie.

Three things follow.

\textbf{The location did not replicate.} Suite 14's 6.6--7.4M window sits squarely inside the
minGRU-lead region on this recipe, with minGRU ahead by 0.13--0.18.

\textbf{There are two crossings, not one.} Attention leads the \emph{first} evaluation (0.836M) on every
seed here --- the opposite of suite 14, where minGRU led from its first logged point. The early
crossing at 1.05M is a phenomenon suite 14 could not have seen: its first evaluation was at 0.82M,
essentially on top of the crossing, and it had one seed.

\textbf{Absolute losses are not comparable across the two boxes.} At \textasciitilde8.2M, attention is 5.434 here
against 5.136 on the 3070 Ti --- a \textasciitilde0.3 nat gap. The curves are not vertically interchangeable, and
no table in this paper stacks them.

\subsection{4.3 The learning-rate horizon is a first-order confound (suites 23 and 24)}\label{the-learning-rate-horizon-is-a-first-order-confound-suites-23-and-24}

Suite 23 asked for an independent 20M-token pair on the same batch and evaluation recipe. It
recovered the early crossing exactly --- \textbf{1.047M}, per-seed 1.03--1.09M --- and moved the late one to
\textbf{14.582M} (per-seed 13.709, 14.490, 14.929, 15.058, 15.193M). At 12.3M, minGRU still led by
\textbf{+0.051} on every seed; at 19.677M attention led by only \textbf{-0.065} {[}-0.081, -0.050{]}, against
suite 22's -0.193 at the same marker.

The cause is mechanical and was verified in the configs rather than inferred: the cosine schedule
followed \texttt{max\_steps} (1220) instead of the 50M horizon (3051). At the 12.3M marker the
20M-horizon learning rate is approximately \textbf{47\%} of the 50M-horizon learning rate at the same
token count. Two runs that are identical in architecture, data, batch, seed set and evaluation
protocol are on materially different learning rates at the token where the ranking flips.

Suite 24 is the corrected prefix: identical to suite 23 except \texttt{lr\_max\_steps\ =\ 3051}.

{%
{\small\begin{longtable}[]{@{}
  >{\raggedright\arraybackslash}p{(\linewidth - 6\tabcolsep) * \real{0.2500}}
  >{\raggedright\arraybackslash}p{(\linewidth - 6\tabcolsep) * \real{0.2500}}
  >{\raggedright\arraybackslash}p{(\linewidth - 6\tabcolsep) * \real{0.2500}}
  >{\raggedright\arraybackslash}p{(\linewidth - 6\tabcolsep) * \real{0.2500}}@{}}
\toprule\noalign{}
\begin{minipage}[b]{\linewidth}\raggedright
Actual tokens
\end{minipage} & \begin{minipage}[b]{\linewidth}\raggedright
Attention {[}95\%{]}
\end{minipage} & \begin{minipage}[b]{\linewidth}\raggedright
minGRU {[}95\%{]}
\end{minipage} & \begin{minipage}[b]{\linewidth}\raggedright
gap A-G {[}95\%{]}
\end{minipage} \\
\midrule\noalign{}
\endhead
\bottomrule\noalign{}
\endlastfoot
0.836M & 6.778 {[}6.765, 6.791{]} & 6.831 {[}6.819, 6.844{]} & \textbf{-0.053} {[}-0.055, -0.051{]} \\
4.112M & 5.888 {[}5.863, 5.914{]} & 5.661 {[}5.637, 5.684{]} & \textbf{+0.227} {[}+0.214, +0.241{]} \\
6.570M & 5.580 {[}5.572, 5.589{]} & 5.402 {[}5.390, 5.414{]} & \textbf{+0.178} {[}+0.168, +0.189{]} \\
7.389M & 5.504 {[}5.489, 5.518{]} & 5.350 {[}5.330, 5.370{]} & \textbf{+0.153} {[}+0.141, +0.166{]} \\
8.208M & 5.432 {[}5.401, 5.464{]} & 5.299 {[}5.267, 5.332{]} & \textbf{+0.133} {[}+0.125, +0.141{]} \\
12.304M & 5.122 {[}5.099, 5.144{]} & 5.120 {[}5.097, 5.144{]} & +0.001 {[}-0.010, +0.013{]} (tie) \\
19.677M & \textbf{4.744} {[}4.734, 4.754{]} & 4.936 {[}4.924, 4.947{]} & \textbf{-0.192} {[}-0.199, -0.185{]} \\
\end{longtable}}
}

Crossings: \textbf{1.038M} and \textbf{12.335M}; per-seed late crossings 11.930, 12.287, 12.365, 12.564,
12.579M. This is an independent ten-job draw that lands on suite 22's crossing to within 0.02M and
reproduces the 19.677M gap to within \textbf{0.001}. Suite 23's 14.6M was the schedule, not seed noise.

This is, to us, the most actionable single finding in the paper. \textbf{Truncating a run for a cheaper
comparison silently retunes the learning rate at every token you are comparing at.} A study that
stops a 50M-token recipe at 20M to save compute is not measuring a prefix of that recipe.

\subsection{4.4 Batch size changes the early ordering (suite 25)}\label{batch-size-changes-the-early-ordering-suite-25}

Suite 22 changed GPU and batch size together. Suite 25 holds batch at suite 14's value of 8 on the
GH200, at suite 14's 8.192M budget, with the cosine horizon matching the run (2000 steps).

{%
{\small\begin{longtable}[]{@{}
  >{\raggedright\arraybackslash}p{(\linewidth - 6\tabcolsep) * \real{0.2500}}
  >{\raggedright\arraybackslash}p{(\linewidth - 6\tabcolsep) * \real{0.2500}}
  >{\raggedright\arraybackslash}p{(\linewidth - 6\tabcolsep) * \real{0.2500}}
  >{\raggedright\arraybackslash}p{(\linewidth - 6\tabcolsep) * \real{0.2500}}@{}}
\toprule\noalign{}
\begin{minipage}[b]{\linewidth}\raggedright
Actual tokens
\end{minipage} & \begin{minipage}[b]{\linewidth}\raggedright
Attention {[}95\%{]}
\end{minipage} & \begin{minipage}[b]{\linewidth}\raggedright
minGRU {[}95\%{]}
\end{minipage} & \begin{minipage}[b]{\linewidth}\raggedright
gap A-G {[}95\%{]}
\end{minipage} \\
\midrule\noalign{}
\endhead
\bottomrule\noalign{}
\endlastfoot
0.823M & 6.736 {[}6.698, 6.774{]} & 6.567 {[}6.521, 6.614{]} & \textbf{+0.168} {[}+0.160, +0.177{]} \\
4.100M & 5.834 {[}5.799, 5.868{]} & 5.721 {[}5.686, 5.755{]} & \textbf{+0.113} {[}+0.102, +0.124{]} \\
6.558M & 5.507 {[}5.474, 5.539{]} & 5.462 {[}5.420, 5.504{]} & \textbf{+0.045} {[}+0.028, +0.062{]} \\
7.377M & 5.425 {[}5.391, 5.459{]} & 5.395 {[}5.359, 5.431{]} & \textbf{+0.030} {[}+0.017, +0.043{]} \\
\end{longtable}}
}

\textbf{No crossing occurs in any logged evaluation, on the mean curves or on any individual seed.}
minGRU leads from the first evaluation on all five seeds and the gap shrinks monotonically ---
0.168 $\rightarrow$ 0.113 $\rightarrow$ 0.045 $\rightarrow$ 0.030 --- without changing sign through 7.38M tokens.

This isolates batch size from hardware. At batch 32 on this GPU, attention leads the first
evaluation on every seed; at batch 8 on the \emph{same} GPU, minGRU leads the first evaluation on
every seed. Suite 25 reproduces suite 14's \emph{early shape}, which suite 22 does not, while failing
to reproduce suite 14's 6.6--7.4M overtake, which suite 22 also does not.

One honest gap: the final 8.192M \texttt{evaluate()} is written to \texttt{done.best\_val} and not logged as an
eval row, so the last \textbf{paired} marker here is 7.377M. \texttt{best\_val} (attention 5.371 {[}5.330,
5.412{]}, minGRU 5.353 {[}5.319, 5.387{]}) is a minimum over all evaluations, not a paired snapshot,
and we do not rank on it. Confidence is \textbf{High} for ``no flip through 7.38M'' and \textbf{Medium} for
any statement about 8.2M.

\subsection{4.5 A matched-batch 50M ranking (suite 26)}\label{a-matched-batch-50m-ranking-suite-26}

Suite 22's ten-arm 50M board mixed batch 32 / \texttt{eval\_iters\ =\ 20} with batch 96 / \texttt{eval\_iters\ =\ 4},
because the token budget was drained by repacking leftover jobs at a wider microbatch. That makes
it a record of what finished, not a bakeoff. Suite 26 reruns the eight drifted arms at batch 32
and \texttt{eval\_iters\ =\ 20} from step 0, 40 jobs, 40/40 completed.

{%
{\small\begin{longtable}[]{@{}
  >{\raggedright\arraybackslash}p{(\linewidth - 8\tabcolsep) * \real{0.2000}}
  >{\raggedright\arraybackslash}p{(\linewidth - 8\tabcolsep) * \real{0.2000}}
  >{\raggedright\arraybackslash}p{(\linewidth - 8\tabcolsep) * \real{0.2000}}
  >{\raggedright\arraybackslash}p{(\linewidth - 8\tabcolsep) * \real{0.2000}}
  >{\raggedright\arraybackslash}p{(\linewidth - 8\tabcolsep) * \real{0.2000}}@{}}
\toprule\noalign{}
\begin{minipage}[b]{\linewidth}\raggedright
Rank
\end{minipage} & \begin{minipage}[b]{\linewidth}\raggedright
Arm
\end{minipage} & \begin{minipage}[b]{\linewidth}\raggedright
Source
\end{minipage} & \begin{minipage}[b]{\linewidth}\raggedright
Mean val {[}95\%{]}
\end{minipage} & \begin{minipage}[b]{\linewidth}\raggedright
Seed range
\end{minipage} \\
\midrule\noalign{}
\endhead
\bottomrule\noalign{}
\endlastfoot
1 & attention & s22 & \textbf{4.222} {[}4.204, 4.240{]} & 4.204--4.237 \\
2 & hybrid\_mingru10\_attn2 & s26 & \textbf{4.232} {[}4.210, 4.254{]} & 4.210--4.249 \\
3 & hybrid\_gdn\_periodic & s26 & 4.290 {[}4.271, 4.309{]} & 4.271--4.308 \\
4 & hybrid\_gdn\_bookend & s26 & 4.301 {[}4.278, 4.324{]} & 4.277--4.319 \\
5 & hybrid\_gdn10\_attn2 & s26 & 4.314 {[}4.292, 4.336{]} & 4.290--4.331 \\
6 & hybrid\_mamba10\_attn2 & s26 & 4.333 {[}4.312, 4.353{]} & 4.313--4.353 \\
7 & gdn & s26 & 4.441 {[}4.419, 4.462{]} & 4.422--4.461 \\
8 & mingru & s22 & 4.449 {[}4.423, 4.476{]} & 4.430--4.473 \\
9 & mamba2 & s26 & 4.596 {[}4.562, 4.630{]} & 4.572--4.633 \\
10 & mla & s26 & 4.627 {[}4.604, 4.650{]} & 4.609--4.647 \\
\end{longtable}}
}

\begin{quote}
\textbf{Horizon and a rejected correction (2026-08-24).} Every row above is measured at the same
evaluation point, \textbf{49,987,584 tokens}: the forty \texttt{crossover50m\_matched32} runs and the
\texttt{crossover50m} attention and minGRU arms all carry that point. Where this section compares
movement against suite 22 below, those suite-22 values are measured at \textbf{49,594,368} tokens
instead --- a 0.8\% shorter horizon --- because four \texttt{crossover50m} arms end on the shorter grid.
That difference is smaller than every rank gap discussed, but it is a cross-horizon comparison
and is labelled as one.

A later rebuild of this manuscript (\S7.4 item 8) replaced row 2 with \textbf{4.275 {[}4.245, 4.305{]}}
and reported its interval as disjoint from attention's, on the grounds that ``the 4.232 figure
appears in no surviving run''. \textbf{We have restored 4.232 and reject that correction.} The figure
is \texttt{hybrid\_mingru10\_attn2} in \texttt{crossover50m\_matched32}, and recomputing it from those five runs
gives 4.2319 {[}4.2099, 4.2539{]} at 49,987,584 tokens --- the same point as the attention row it is
compared against. The replacement 4.275 is the same arm in \texttt{crossover50m}, where its seeds end
at 49,594,368. The correction therefore swapped a matched-horizon comparison for a mismatched
one, in the course of documenting horizon mismatch as a defect.

The genuine weakness of this row is the one we already state: it is a \textbf{cross-suite} comparison,
because suite 26 never reran attention or minGRU at 50M. That is gap E2, and it caps this board
at Medium-High confidence until those ten jobs are run.
\end{quote}

Two arms moved substantially relative to the mixed-recipe board, and both are arms that had
trained at batch 96: \texttt{hybrid\_mamba10\_attn2} improved from 4.604 to \textbf{4.333} (-0.271) and
\texttt{hybrid\_mingru10\_attn2} from 4.275 to \textbf{4.232}. GDN-family arms were already at batch 32 and moved
only at seed scale (e.g.~\texttt{hybrid\_gdn\_periodic} 4.280 $\rightarrow$ 4.290).

The matched conclusions are narrower than the drifted board suggested:

\begin{itemize}
\tightlist
\item
  Attention has the \textbf{best mean}, but is \textbf{not separable} from a 10$\times$minGRU + 2$\times$attention hybrid
  at n = 5; the intervals overlap.
\item
  Adding two attention layers to a recurrent stack is worth roughly \textbf{0.22} nats over the pure
  recurrent stack (minGRU 4.449 $\rightarrow$ hybrid 4.232) and roughly \textbf{0.26} over pure Mamba-2
  (4.596 $\rightarrow$ 4.333). This is the largest single effect on the board.
\item
  Pure Mamba-2 and MLA are last, clearly separated from everything above them.
\item
  Nothing here is a throughput claim. Quality ranking and tokens-per-second are separate outcomes
  and are never combined in this paper.
\end{itemize}

Confidence is \textbf{High} for the eight rerun arms and \textbf{Medium-High} for the combined ranking,
because the attention and minGRU cells are suite 22's sample rather than a fresh 50M draw.

\subsection{4.6 Summary: what moved the crossing}\label{summary-what-moved-the-crossing}

Every row below is n = 5 except the first.

{%
{\small\begin{longtable}[]{@{}
  >{\raggedright\arraybackslash}p{(\linewidth - 6\tabcolsep) * \real{0.2500}}
  >{\raggedright\arraybackslash}p{(\linewidth - 6\tabcolsep) * \real{0.2500}}
  >{\raggedright\arraybackslash}p{(\linewidth - 6\tabcolsep) * \real{0.2500}}
  >{\raggedright\arraybackslash}p{(\linewidth - 6\tabcolsep) * \real{0.2500}}@{}}
\toprule\noalign{}
\begin{minipage}[b]{\linewidth}\raggedright
Configuration
\end{minipage} & \begin{minipage}[b]{\linewidth}\raggedright
Suite
\end{minipage} & \begin{minipage}[b]{\linewidth}\raggedright
Early crossing
\end{minipage} & \begin{minipage}[b]{\linewidth}\raggedright
Late crossing
\end{minipage} \\
\midrule\noalign{}
\endhead
\bottomrule\noalign{}
\endlastfoot
RTX 3070 Ti, bs8, 8.192M budget, \textbf{1 seed} & 14 & not observed (minGRU leads eval 1) & \textbf{6.6--7.4M} \\
GH200, bs32, 50M cosine, 50M budget & 22 & \textbf{1.049M} & \textbf{12.353M} \\
GH200, bs32, 50M cosine, 20M budget (independent draw) & 24 & \textbf{1.038M} & \textbf{12.335M} \\
GH200, bs32, \textbf{20M cosine}, 20M budget & 23 & \textbf{1.047M} & \textbf{14.582M} \\
GH200, \textbf{bs8}, own cosine, 8.192M budget & 25 & not observed (minGRU leads eval 1) & \textbf{none through 7.38M} \\
\end{longtable}}
}

Reading down the table: the early crossing is stable at \textasciitilde1.04--1.05M across every batch-32 recipe
and absent at batch 8. The late crossing moves by 2.248M tokens on a change to the learning-rate
horizon alone, and does not occur at all within the original budget at the original batch size.
The one quantity that was reported as the finding --- the token of the flip --- is the one quantity
that is not stable.

\begin{center}\rule{0.5\linewidth}{0.5pt}\end{center}

\section{5. The same phenomenon on the optimizer axis}\label{the-same-phenomenon-on-the-optimizer-axis}

The mixer study varied architecture with everything else fixed. The obvious next axis is the
optimizer, and it is a well-posed target precisely because optimizer papers are so often
introduced on short curves.

\subsection{5.1 Funnel design}\label{funnel-design}

Sixteen candidates, one engine, one equal-token protocol, on an Apple M5 Pro. Fourteen are
implemented natively in Rust and Metal with parity fixtures: Muon NS5, Muon NS3, Polar Express
Muon, NorMuon, Muown, MONA, AdamW, Cautious AdamW, Lion, Cautious Lion, momentum SGD, Sophia,
Schedule-Free AdamW, and Prodigy.

Two are recorded as \textbf{systems-blocked exclusions} rather than dropped: MiMuon (\texttt{mimuon\_adamw}) requires exact
singular-gap routing, which requires a per-matrix SVD; SOAP (\texttt{soap\_adamw}) requires a periodic symmetric
eigendecomposition. Metal/MPS exposes neither, and an Accelerate fallback would force a host
synchronization on every refresh, converting the arm into a measurement of the fallback. They are
not silent substitutions and they are \textbf{not evidence about those methods}.

The funnel is five stages on an atomic, resumable job ledger, written down before the runs:

\begin{enumerate}
\def\labelenumi{\arabic{enumi}.}
\tightlist
\item
  Five-point logarithmic learning-rate sweep at 16M parameters, 100 steps, seed 1337.
\item
  Every stable LR winner for 500 steps, seed 1337.
\item
  The best of those, plus mandatory AdamW and Muon NS5 anchors, for 1,000 steps at seeds 42 and 2026.
\item
  The top four at the \texttt{exact\_128m} configuration (128,367,988 parameters) for 500 steps, seed 1337.
\item
  The top two at 128M for 1,000 steps, seeds 42 and 2026.
\end{enumerate}

Selection is by mean equal-token validation bits-per-byte, with time-to-loss, memory footprint
and step time as tie-breakers \emph{only} when intervals overlap.

\subsection{5.2 The reversal}\label{the-reversal}

Stage 3 --- 16M parameters, 1,000 steps, two seeds. We report both seeds and two intervals: the
$\pm$1.96$\cdot$SE interval originally computed, and the correct Student-\emph{t} interval for n = 2 (t$_1$ = 12.706).

{%
{\small\begin{longtable}[]{@{}
  >{\raggedright\arraybackslash}p{(\linewidth - 10\tabcolsep) * \real{0.1667}}
  >{\raggedright\arraybackslash}p{(\linewidth - 10\tabcolsep) * \real{0.1667}}
  >{\raggedright\arraybackslash}p{(\linewidth - 10\tabcolsep) * \real{0.1667}}
  >{\raggedright\arraybackslash}p{(\linewidth - 10\tabcolsep) * \real{0.1667}}
  >{\raggedright\arraybackslash}p{(\linewidth - 10\tabcolsep) * \real{0.1667}}
  >{\raggedright\arraybackslash}p{(\linewidth - 10\tabcolsep) * \real{0.1667}}@{}}
\toprule\noalign{}
\begin{minipage}[b]{\linewidth}\raggedright
Candidate
\end{minipage} & \begin{minipage}[b]{\linewidth}\raggedright
Seed 42
\end{minipage} & \begin{minipage}[b]{\linewidth}\raggedright
Seed 2026
\end{minipage} & \begin{minipage}[b]{\linewidth}\raggedright
Mean BPB
\end{minipage} & \begin{minipage}[b]{\linewidth}\raggedright
$\pm$1.96$\cdot$SE
\end{minipage} & \begin{minipage}[b]{\linewidth}\raggedright
95\% \emph{t}$_1$ interval
\end{minipage} \\
\midrule\noalign{}
\endhead
\bottomrule\noalign{}
\endlastfoot
MONA & 2.183040 & 2.178756 & \textbf{2.180898} & 0.004198 & {[}2.1537, 2.2081{]} \\
Polar Muon & 2.179047 & 2.185596 & \textbf{2.182322} & 0.006418 & {[}2.1407, 2.2239{]} \\
NorMuon & 2.236222 & 2.228217 & 2.232220 & 0.007845 & {[}2.1814, 2.2831{]} \\
Muon NS5 & 2.225931 & 2.239343 & 2.232637 & 0.013144 & {[}2.1474, 2.3178{]} \\
Muon NS3 & 2.262745 & 2.255586 & 2.259166 & 0.007016 & {[}2.2137, 2.3046{]} \\
Muown & 2.583694 & 2.585490 & 2.584592 & 0.001760 & {[}2.5732, 2.5960{]} \\
AdamW (control) & 11.378829 & 11.179040 & 11.278935 & 0.195790 & {[}10.010, 12.548{]} \\
\end{longtable}}
}

The same four candidates at the \texttt{exact\_128m} configuration (128,367,988 parameters), 500 steps, seed 1337 (n = 1, no
interval):

{%
{\small\begin{longtable}[]{@{}ll@{}}
\toprule\noalign{}
Candidate & Val BPB \\
\midrule\noalign{}
\endhead
\bottomrule\noalign{}
\endlastfoot
Polar Muon & \textbf{2.420958} \\
NorMuon & 2.567474 \\
Muon NS5 & 2.601096 \\
MONA & \textbf{2.950696} \\
\end{longtable}}
}

\textbf{MONA goes from first to last}, by more than half a bit per byte, on the same data, the same
schedule and the same tuned learning rate that won at 16M.

Stage 5, 128M, 1,000 steps, seeds 42 and 2026:

{%
{\small\begin{longtable}[]{@{}lllll@{}}
\toprule\noalign{}
Candidate & Seed 42 & Seed 2026 & Mean BPB & 95\% \emph{t}$_1$ interval \\
\midrule\noalign{}
\endhead
\bottomrule\noalign{}
\endlastfoot
Polar Muon & 2.167282 & 2.172555 & \textbf{2.169919} & {[}2.1364, 2.2034{]} \\
NorMuon & 2.202106 & 2.200184 & 2.201145 & {[}2.1889, 2.2134{]} \\
\end{longtable}}
}

Polar Express Muon is the champion. The locked winner passed an exact checkpoint-and-replay gate:
loss delta 1.43e-6, gradient-norm delta 0.0, sampled master-weight max delta 2.56e-6, 1,707
dispatches, 13,483 MB resident, on a model of exactly \textbf{128,367,988} parameters, with no swap
pressure and no recorded failures.

\subsection{5.3 A correction to our own statistics}\label{a-correction-to-our-own-statistics}

Our earlier write-up of this funnel stated that at stage 5 ``the intervals do not overlap, so the
champion is decided without ever reaching for a systems tie-breaker.'' That was computed with a
normal (z = 1.96) multiplier on a two-seed standard error. With the correct t$_1$ = 12.706 multiplier,
Polar's interval {[}2.1364, 2.2034{]} and NorMuon's {[}2.1889, 2.2134{]} \textbf{do overlap}.

What survives is weaker but still real: \textbf{both} Polar seeds (2.167282, 2.172555) are below
\textbf{both} NorMuon seeds (2.200184, 2.202106), a sign-consistent 2-of-2 ordering with a mean
separation of 0.031 BPB. The correct statement is that the selection rests on a consistent
ordering at n = 2, not on separated 95\% intervals. The same correction applies at stage 3, where
under t$_1$ intervals the top four candidates are \textbf{mutually inseparable} --- which makes the
subsequent 128M reversal less surprising and more damning of the screen.

The defect was in the funnel tooling, not only in the write-up, and it had teeth: the interval
also gated \texttt{confidence\_interval\_overlaps\_best}, which decides whether the declared systems
tie-breakers (time-to-loss, footprint, step time) are invoked at all. A single-seed arm was
additionally recorded with a \textbf{zero-width} interval --- infinite precision for an arm measured
once --- so it could be declared separated from every rival. \path{nanolab/native_funnel.py} now uses a
Student-t multiplier with df = n - 1, reports an \textbf{undefined} (infinite) interval at n = 1, and
\textbf{skips} the systems tie-breakers entirely when any arm has fewer than three seeds, recording
\texttt{confidence\_interval\_overlaps\_best\ =\ None} rather than \texttt{False} so that a check which could not run
is not recorded as a check that ran. Stored rankings were recomputed with the original values
preserved under \texttt{ci95\_legacy\_z}. The recorded champion is unchanged. Full detail in
\path{docs/ISSUES_AND_GAPS_2026-08-22.md} \S2.3.

\subsection{5.4 The control that failed loudly}\label{the-control-that-failed-loudly}

AdamW appears at 11.28 BPB in the stage-3 table. It is worth explaining rather than hiding.

At the 100-step LR sweep, AdamW at lr 0.0012 finished at 3.593993 --- behind the five Muon-family
arms (3.353--3.376) but ahead of Muown (3.642), Sophia (3.643), Cautious Lion (3.739), momentum SGD
(3.754), Lion (3.755) and Schedule-Free AdamW (3.778). It looked like a healthy control. At 500 steps it was 8.489102. At 1,000 steps across two seeds it was 11.278935.

Nothing broke: no NaN, no divergence, finite telemetry throughout. The learning rate that looked
best over 100 steps simply does not survive 1,000, and short-horizon EMA lag hid that completely
at the point where the decision was made. The narrow, transferable version: \textbf{a 100-step
learning-rate selection is not a learning-rate selection.} It is a measurement of which method
has the least EMA lag at 100 steps.

This does not overturn independent evidence that AdamW is quality-competitive. A separate
one-seed, 4.096M-token, width-768 bake-off on the 3070 Ti ranked Lion first at 5.557 and AdamW
second at 5.565, ahead of Muon at 5.643 --- which is exactly why AdamW is a \emph{mandatory} anchor in
this funnel rather than an optional arm. In the same bake-off Prodigy soft-diverged to 20.16 and
Schedule-Free AdamW plateaued at 7.918 at their defaults. In the native funnel Prodigy is a
recorded numerical exclusion rather than a missing result: two of its five learning-rate sweep
points terminated non-finite after two logged steps, its best finite screen result was 11.889982
BPB at lr 0.25 --- an order of magnitude behind every other candidate --- and its stage-2 500-step job
at that same learning rate also terminated non-finite.

\subsection{5.5 The mirror image: downward transfer fails too}\label{the-mirror-image-downward-transfer-fails-too}

If the story were ``small-scale rankings underestimate the eventual winner,'' it would be a tidy
lesson and half true. We ran the 128M champion back down against Muon NS5 on the small \texttt{sota}
model, 3,000 steps, two seeds, at matched throughput on the quality-oriented flash path:

{%
{\small\begin{longtable}[]{@{}llll@{}}
\toprule\noalign{}
Seed & NS5 & Polar & $\Delta$ \\
\midrule\noalign{}
\endhead
\bottomrule\noalign{}
\endlastfoot
1337 & \textbf{2.122381} & 2.127708 & +0.005327 \\
42 & \textbf{2.127647} & 2.131005 & +0.003358 \\
mean & \textbf{2.125014} & 2.129357 & \textbf{+0.004342} \\
\end{longtable}}
}

The 128M champion loses at small scale, on both seeds. The champion's tuned learning rate of 0.05
applied at the small scale was worse still. So the ranking is not a biased-but-monotone estimate of
the large-scale ranking; it is \textbf{a property of the scale}, with no safe direction of
extrapolation. In this funnel the eventual champion was third at the LR sweep and second at 1,000
steps --- a shortlist of two would have discarded it at stage 1 had the LR screen been the only cut.

\begin{center}\rule{0.5\linewidth}{0.5pt}\end{center}

\section{6. The pattern is not confined to these two axes}\label{the-pattern-is-not-confined-to-these-two-axes}

\S4 and \S5 each show one axis on which a ranking moves with the measurement recipe. If those were
the only two, the honest reading would be that we found two awkward special cases. They are not.
Assembling the rest of our own record --- suites that predate this paper, and a second hardware
backend --- turns up the same structure on six further axes. We set the catalogue out here because a
reader is entitled to ask how often this happens before accepting \S8's recommendations, and
because several of these results were sitting unreported in our own artifacts until we audited
them for this draft.

The evidence in this section is \textbf{weaker than in \S4}, and uniformly so: most of it is n = 1 at
seed 1337 on a single 8 GB laptop GPU, with no confidence intervals anywhere. We label each entry
accordingly and we draw no quantitative conclusion from any of them. The claim this section
supports is qualitative and, we think, the more important one: \emph{rank reversal under a change of
measurement recipe is the common case, not a curiosity of the attention/minGRU pair.}

\subsection{6.1 Token budget alone reorders a five-arm board}\label{token-budget-alone-reorders-a-five-arm-board}

Suite 13 ranked five mixers at 2.048M tokens; suite 14 extended three of them to 8.192M under a
byte-identical configuration --- same 12L/768d shape, block 512, batch 8, Muon at 6e-4, FineWeb-edu,
seed 1337, \texttt{mixer\_chunk\ =\ 32}. Only the budget changed.

{%
{\small\begin{longtable}[]{@{}lll@{}}
\toprule\noalign{}
arm & 2.048M (suite 13) & 8.192M (suite 14) \\
\midrule\noalign{}
\endhead
\bottomrule\noalign{}
\endlastfoot
minGRU & \textbf{5.8374} (1st of 5) & 5.1545 (2nd of 3) \\
Gated DeltaNet & 5.9939 (2nd) & --- \\
Mamba-2 & 6.0404 (3rd) & 5.3831 (3rd of 3) \\
attention & 6.0733 (\textbf{4th}) & \textbf{5.1362} (\textbf{1st}) \\
MLA & 6.1556 (5th) & --- \\
\end{longtable}}
}

Attention --- the default the other four are proposed as alternatives to --- is fourth of five at the
short budget and first at the long one. This is the cleanest statement of the paper's thesis we
have, because nothing but the number of tokens differs, and it is a \emph{board} reordering rather than
a single crossing. It is also n = 1 per arm, and the 8.192M values are \texttt{best\_val} fields rather
than paired snapshots, so it should be read as a demonstration of the phenomenon and not as a
measurement of its size. \S4's n = 5 replication exists precisely because this one does not carry
an interval.

\subsection{6.2 The choice of metric reorders the same arms}\label{the-choice-of-metric-reorders-the-same-arms}

The five mixers of suite 13 were separately benchmarked for throughput at fixed shape (suite 17,
batch 8, context 512, 124M parameters, same GPU):

{%
{\small\begin{longtable}[]{@{}
  >{\raggedright\arraybackslash}p{(\linewidth - 6\tabcolsep) * \real{0.2500}}
  >{\raggedright\arraybackslash}p{(\linewidth - 6\tabcolsep) * \real{0.2500}}
  >{\raggedright\arraybackslash}p{(\linewidth - 6\tabcolsep) * \real{0.2500}}
  >{\raggedright\arraybackslash}p{(\linewidth - 6\tabcolsep) * \real{0.2500}}@{}}
\toprule\noalign{}
\begin{minipage}[b]{\linewidth}\raggedright
arm
\end{minipage} & \begin{minipage}[b]{\linewidth}\raggedright
quality rank (suite 13)
\end{minipage} & \begin{minipage}[b]{\linewidth}\raggedright
throughput (suite 17)
\end{minipage} & \begin{minipage}[b]{\linewidth}\raggedright
throughput rank
\end{minipage} \\
\midrule\noalign{}
\endhead
\bottomrule\noalign{}
\endlastfoot
MLA & \textbf{5th} (6.1556, worst) & 9,334.57 tok/s & \textbf{1st} \\
attention & 4th (6.0733) & 7,923.37 tok/s & 2nd \\
minGRU & \textbf{1st} (5.8374, best) & 6,674.23 tok/s & 3rd \\
Mamba-2 & 3rd (6.0404) & 333.23 tok/s & 4th \\
Gated DeltaNet & 2nd (5.9939) & 238.18 tok/s & 5th \\
\end{longtable}}
}

MLA is last on quality and first on throughput; minGRU is first on quality and third on
throughput. The two orderings are close to inverted at the top. Neither metric is wrong --- they
answer different questions --- but a screen that reports one and is read as if it ranked the other
will select the opposite arm. Note also the scale of the throughput spread: Mamba-2 and Gated
DeltaNet are 24$\times$ and 33$\times$ slower than attention here, which is a difference in kind rather than
degree and is what motivated \S6.5.

\subsection{6.3 Within a single suite, speed rank and quality rank invert}\label{within-a-single-suite-speed-rank-and-quality-rank-invert}

Suite 18 stacked memory-residency optimizations on a 124M attention model and ranked the arms by
throughput. Its \texttt{best\_val} figures run the other way:

{%
{\small\begin{longtable}[]{@{}llll@{}}
\toprule\noalign{}
arm & peak tok/s & MFU & \texttt{best\_val} \\
\midrule\noalign{}
\endhead
\bottomrule\noalign{}
\endlastfoot
\texttt{gpu\_opt\_bs32} & \textbf{13.7K} (1st) & 25.5\% & \textbf{4.9065} (3rd) \\
\texttt{gpu\_opt\_validate} & 13.6K (2nd) & 25.3\% & 4.8084 (2nd) \\
\texttt{gpu\_max} & 11.9K (3rd) & 22.1\% & \textbf{4.8001} (1st) \\
\end{longtable}}
}

The ordering is exactly reversed. The three arms differ in batch size (16 / 24 / 32) at a fixed
step count, so they do not see equal token budgets and this is emphatically \emph{not} a controlled
quality comparison --- which is the point. The suite was built to answer a systems question, its
throughput conclusion is sound, and reading its ladder as a recipe ladder selects the worst of the
three models. We flag this one with some embarrassment: \texttt{gpu\_opt\_bs32}'s \texttt{best\_val} was the single
cell absent from that suite's own results table until this audit, and its absence is what made the
inversion invisible for two months.

\subsection{6.4 Batch size reorders throughput among mixers}\label{batch-size-reorders-throughput-among-mixers}

Suites 22 and 26 ran overlapping hybrid arms at batch 96 and batch 32 respectively on the GH200,
with \texttt{mixer\_chunk\ =\ 32} in both. Median per-step throughput across five seeds:

{%
{\small\begin{longtable}[]{@{}lll@{}}
\toprule\noalign{}
arm & batch 96 & batch 32 \\
\midrule\noalign{}
\endhead
\bottomrule\noalign{}
\endlastfoot
\texttt{hybrid\_mamba10\_attn2} & \textbf{43,941 tok/s} & 17,753 tok/s \\
\texttt{hybrid\_gdn10\_attn2} & 31,333 tok/s & \textbf{24,279 tok/s} \\
\end{longtable}}
}

At batch 96 the Mamba-2 hybrid is \textasciitilde1.4$\times$ faster than the GDN hybrid, with disjoint seed ranges
(43,913--43,946 against 31,330--31,383). At batch 32 the ordering reverses. Mamba-2's chunked scan
needs batch parallelism to fill the device; the delta-rule kernel degrades more gracefully. Two
caveats matter. First, the quality comparison between these two suites is confounded --- batch size
and \texttt{lr\_max\_steps} both change --- but the \emph{throughput} comparison is not, because the learning-rate
schedule cannot affect tokens per second. Second, these suites ran two jobs per GPU (measured
concurrency 1.97$\times$), so the absolute figures are contended and one \texttt{mamba2} seed spans
15,493--38,365 tok/s. We therefore report the reversal and not the magnitudes, and a single-tenant
re-measurement is on the list in \S8.4.

\subsection{6.5 The kernel implementation decides feasibility, not just speed}\label{the-kernel-implementation-decides-feasibility-not-just-speed}

Suite 19 replaced the O(T) sequential scans for Mamba-2 and Gated DeltaNet with chunk-parallel
kernels, verified against the sequential reference for output and input-gradient parity at 1e-5
with fp32 accumulation and non-divisible-length coverage:

{%
{\small\begin{longtable}[]{@{}
  >{\raggedright\arraybackslash}p{(\linewidth - 4\tabcolsep) * \real{0.3333}}
  >{\raggedright\arraybackslash}p{(\linewidth - 4\tabcolsep) * \real{0.3333}}
  >{\raggedright\arraybackslash}p{(\linewidth - 4\tabcolsep) * \real{0.3333}}@{}}
\toprule\noalign{}
\begin{minipage}[b]{\linewidth}\raggedright
arm
\end{minipage} & \begin{minipage}[b]{\linewidth}\raggedright
before
\end{minipage} & \begin{minipage}[b]{\linewidth}\raggedright
after
\end{minipage} \\
\midrule\noalign{}
\endhead
\bottomrule\noalign{}
\endlastfoot
Mamba-2 @ bs8/ctx512 & 333 tok/s & 3,224 tok/s (9.7$\times$) \\
Gated DeltaNet @ bs8/ctx512 & 238 tok/s & 482 tok/s, later \textasciitilde1,600 tok/s vectorized \\
Gated DeltaNet @ bs16/ctx1024 & \textbf{OOM} & \textbf{1,100 tok/s @ 4.0 GB} \\
\end{longtable}}
}

The last row is the one that matters for ranking: a configuration that was \emph{unrunnable} becomes
competitive, purely from an implementation change with no effect on the mathematics. Any board
compiled before this work would have recorded GDN at that shape as an absence rather than a
number. Kernel maturity is a property of the software, not the architecture, yet it silently
determines which arms appear on the board at all.

This section also carries the paper's most consequential bug: the first version of the chunked
scan accumulated in bf16 under autocast, and CPU-only tests did not catch the drift. Every SSM
quality number downstream would have been wrong. The fix is to force fp32 inside the scan.

\subsection{6.6 The hardware backend reorders both quality and throughput}\label{the-hardware-backend-reorders-both-quality-and-throughput}

We maintain a second, independent implementation of the same architecture in Rust and Metal on an
M5 Pro. Comparing it against the CUDA reference at the same \texttt{sota} shape gives a quality crossover
in the training horizon:

{%
{\small\begin{longtable}[]{@{}llll@{}}
\toprule\noalign{}
horizon & Metal-native (f32) & 3070 Ti CUDA (bf16) & leader \\
\midrule\noalign{}
\endhead
\bottomrule\noalign{}
\endlastfoot
3k steps & 2.0222 & \textbf{1.9944} & CUDA \\
20k, no warmdown & \textbf{1.9178} & 1.9944 & Metal \\
20k + warmdown, seed 1337 & \textbf{1.8969} & 1.9944 & Metal \\
20k + warmdown, seed 42 & \textbf{1.8876} & \textasciitilde1.9860 & Metal \\
100k WSD & \textbf{1.8828} & --- & --- \\
\end{longtable}}
}

The backends swap places somewhere between 3k and 20k steps. This comparison is confounded in
several ways we cannot remove --- different precision (f32 vs bf16), different flash-attention
implementations, and no matched CUDA arm at 20k or 100k --- so we present it as a demonstration that
``which implementation is better'' is horizon-dependent, not as a measurement of either backend.

The backend also reorders the mixers themselves. On CUDA at 124M (suite 17) attention is 1.19$\times$
faster than minGRU; on Metal at \texttt{sota} scale, measured over 40 FineWeb steps with real kernels,
minGRU is \textbf{2.17$\times$ faster than attention} (144,367 vs 66,382 tok/s), and Mamba-2 closes from 24$\times$
slower than attention to 2.0$\times$ slower (33,321 tok/s). The mixer throughput ranking is a property of
the kernel library available on the device.

Most striking, the Metal track independently reproduces \S4's headline shape. At 3k steps with
matched seeded initialization, minGRU leads attention 2.085016 to 2.123010 (and 2.076374 on a
second seed); at 20k steps attention leads 1.89688 to 1.993295. Bias early, capacity late --- the
same direction as \S4, on a different backend, precision, kernel stack and data pipeline. We
report this as suggestive rather than as a replication, because the 20k attention arm uses
\path{-\/-golden-init} while the minGRU arm is seeded (the golden banks contain no \texttt{mingru\_*} tensors),
so the crossing itself is confounded with initialization. One seeded attention run at 20k would
close that gap, and it is the cheapest experiment named anywhere in this paper.

\subsection{6.7 What this catalogue is and is not}\label{what-this-catalogue-is-and-is-not}

Counting \S4, \S5 and this section, we have observed a ranking change places under a change of:
token budget, learning-rate horizon, batch size, model scale, training horizon at fixed scale,
learning rate, evaluation metric, kernel implementation, and hardware backend. In no case did the
architectures, the optimizers or the data change.

We are \textbf{not} claiming that all rankings are arbitrary, that these effects are of comparable size,
or that any particular arm is secretly better. Several entries above are single-seed observations
on one laptop GPU and would not survive a careful replication; \S4 exists because we know the
difference. What the catalogue establishes is narrower and sufficient for \S8: the recipe-dependence
in \S4 is not an anomaly of one architecture pair, so a screening protocol that does not pin and
report its recipe is not measuring a property of the methods it claims to rank. The practices in
\S8 follow from the pattern, not from any single row above.

\section{7. Threats to validity and limitations}\label{threats-to-validity-and-limitations}

\subsection{7.1 What these results do not establish}\label{what-these-results-do-not-establish}

\begin{itemize}
\tightlist
\item
  \textbf{No universal crossing token.} The claim is the opposite. Nothing here supports ``attention
  overtakes minGRU at \emph{X} tokens'' for any \emph{X} outside the exact recipe measured.
\item
  \textbf{No 3070 Ti vs GH200 quality ranking.} The two boxes differ in batch size, evaluation count
  and kernel versions simultaneously. Absolute losses differ by \textasciitilde0.18--0.3 nats at matched token
  markers and the tables are never stacked.
\item
  \textbf{No throughput claim.} Every number in \S4 is held-out cross-entropy. Quality and
  tokens-per-second are separate outcomes.
\item
  \textbf{No mechanism.} ``Bias early, capacity late'' fits the curves and does not predict the crossing
  token, which is what moved. It remains an explanation, not a theorem.
\item
  \textbf{Suite 26 does not rerun attention or minGRU.} Its top two rows are suite 22's sample.
\item
  \textbf{Suite 25 cannot locate a late crossing.} Its horizon ends before any batch-32 recipe crosses.
\end{itemize}

\subsection{7.2 Known limitations}\label{known-limitations}

\begin{itemize}
\tightlist
\item
  \textbf{Single data distribution, single model shape.} Everything is FineWeb-Edu at 12L/768d,
  sequence 512. Whether the schedule sensitivity in \S4.3 persists at other widths, depths or
  sequence lengths is untested.
\item
  \textbf{No $\mu$P, and a single global learning rate.} Runs use standard parametrization with Muon at a
  fixed 6e-4 across every arm, batch size and schedule. This is the largest uncontrolled factor in
  the paper and the one most capable of offering a competing explanation for \S4.3 and \S4.4; it is
  treated at length in \textbf{\S8}, together with the 2 $\times$ 2 follow-up that would settle it and the
  pre-registered reading of each outcome.
\item
  \textbf{Two seeds on the optimizer axis.} \S5's stage-3 and stage-5 intervals are two-seed intervals
  and are, correctly computed, wide. The rank \emph{reversal} is robust (a 0.53 BPB swing at 128M is
  far outside any plausible two-seed noise band); the \emph{ordering within the top four} is not.
\item
  \textbf{\texttt{compile\ =\ False} throughout suites 22--26.} Forced by an Inductor stall on GH200 aarch64.
  Uniform across arms, but a compiled recipe is a different recipe and might move the crossings
  again.
\item
  \textbf{Recipe drift in suite 22's ten-arm board.} Documented in \S4.5 and repaired by suite 26; the
  original board should be read only as ``what finished.''
\item
  \textbf{A kernel bug mid-grid.} Gated DeltaNet produced NaN gradients at finite loss (batch 96
  step 39; batch 32 step 55) because \texttt{gdn\_chunked} exponentiated the full C$\times$C log-ratio grid and
  \texttt{j\ \textgreater{}\ t} entries overflowed, with \texttt{.tril()} masking it in the forward pass. Fixed by taking the
  lower triangle before the exponential, clamping log-ratios at max 0, and clamping alpha. All GDN
  rows in this paper are post-fix. This was a kernel defect, not a mixer-quality result.
\item
  \textbf{A number we reconciled, and the wrong reconciliation we tried first.} The champion run's
  final EMA sliding BPB is \textbf{2.010659} (seed 1337) in \path{research/champion-run.json}, while other
  workspace documents recorded \textbf{2.015756}. We first concluded that 2.015756 was a \texttt{57}$\leftrightarrow$\texttt{75}
  transposition of the audit-7 artifact's 2.015576, on the reasoning that it matched no artifact
  on disk. That was wrong, and acting on it would have deleted a real measurement from four
  documents. There were \textbf{three} champion runs, not two: \texttt{DECISIONS.md} M15, written 2026-07-19
  before Audit 7 existed, independently records the champion at \textbf{2.0158} --- which 2.015756 rounds
  to and 2.015576 does not. 2.015756 is the pre-Audit-7 run's own value; its artifact was
  superseded and deleted, which makes the figure \emph{unverifiable}, not mistyped. ``Matches no
  artifact on disk'' and ``was fabricated'' are different claims, and we conflated them. The citable
  value is unchanged at 2.010659, and no comparison in this paper uses either figure.
  \texttt{champion-run.json} additionally carries \texttt{"locked":\ false} --- see \S8.3 for why it stays that way.
\end{itemize}

\subsection{7.3 Claims withdrawn}\label{claims-withdrawn}

Stating these explicitly is part of the result.

\begin{enumerate}
\def\labelenumi{\arabic{enumi}.}
\tightlist
\item
  \textbf{``The \textasciitilde7M attention/minGRU crossover, replicated.''} Withdrawn as a general claim. Suite 14
  stands as a valid RTX 3070 Ti / batch 8 / one-seed observation and nothing more.
\item
  \textbf{``Metal trainer BPB 2.038 vs CUDA 1.994 confirms numerical correctness.''} Withdrawn. The two
  figures come from different flash implementations and the comparison to the 1.9944 CUDA
  reference is cross-scale, not matched-shape; the same-shape CUDA 128M reference is still
  unrun. Numerical correctness of the Metal trainer rests on what was actually done --- exact
  checkpoint-and-replay to 1e-6 (\S5.2), 57/57 \path{cargo\ test\ -\/-release\ -\/-lib} gates covering the
  native optimizer, checkpoint, telemetry and NaN paths, and an independent MLX oracle for the
  manual GEMM attention backward --- not on that comparison. The claim was attached to the wrong
  evidence. (An ``29/29 oracle tests'' figure circulated in earlier internal write-ups. A previous
  version of this section stated that we could not locate its source; that was wrong. It is
  recorded at \path{experiment-notes/arch-metal/51-m5-128m-optimizer-funnel-preflight.md}, line 291:
  ``Python correctness/orchestration oracle: \textbf{29/29} tests pass (includes induced-swap unlock
  rejection).'' It is a Python orchestration oracle, not a numerical-parity gate, and it is
  distinct from the 57/57 \texttt{cargo} gates above.)
\item
  \textbf{Mixer throughput of \textasciitilde238,000 tok/s (Mamba-2) and \textasciitilde240,000 tok/s (minGRU).} Withdrawn as
  measurement artifacts: the training graph ran the MLP and stem while the mixer path was a
  zero-tensor stub with flash attention skipped, so the figure measured everything except the
  thing being compared. Re-baselined honestly on FineWeb at sequence length 256, the same three
  arms are \textbf{144,367 tok/s} (minGRU), \textbf{66,382} (attention) and \textbf{33,321} (Mamba-2). minGRU is
  genuinely fast; it is not 240,000-fast, and Mamba-2 at this sequence length is \emph{slower} than
  attention rather than several times faster.
\item
  \textbf{``Stage-5 intervals do not overlap.''} Withdrawn and replaced by \S5.3.
\item
  \textbf{A documented \textasciitilde4$\times$ headroom claim in the quantized GEMV inference kernels.} Withdrawn. Measured
  properly, those kernels already run at 62--100\% of the \textasciitilde273 GB/s memory-bandwidth ceiling; the
  4$\times$ came from a microbenchmark that re-uploaded weights on every call, measuring an upload that
  does not happen in production. About 77\% of a quiet decode token is dispatch and barrier
  overhead, not arithmetic. The kernels were not the bottleneck; the number of kernel launches
  was.
\end{enumerate}

\subsection{7.4 Corrections found by auditing our own record}\label{corrections-found-by-auditing-our-own-record}

The claims in \S7.3 were withdrawn as we discovered them, over months. Preparing this draft we did
something different and more uncomfortable: a systematic audit of every number in our own suite
notes, lock files and ledgers against the raw run directories they were derived from. It found
nine defects. We record them here rather than silently fixing them, because a paper arguing that
measurement protocol determines conclusions has no standing to hide its own protocol failures ---
and because the \emph{kind} of error is instructive. None of these was a bad measurement. Every one was
a correct measurement that was then mis-recorded, partially reported, or lost.

\begin{enumerate}
\def\labelenumi{\arabic{enumi}.}
\item
  \textbf{A results table omitted the one cell that contradicted its ranking.} Suite 18 ranked three
  arms by throughput and reported \texttt{best\_val} for two of them. The missing third,
  \texttt{gpu\_opt\_bs32} at \textbf{4.9065}, is the fastest arm and the worst model, and its absence is the
  only reason the speed/quality inversion in \S6.3 went unnoticed for two months.
\item
  \textbf{A headline ranking depended on undocumented restarts.} \texttt{run128m\_20k} is not one run: its
  \texttt{metrics.jsonl} holds \textbf{nine \texttt{start} and eight \texttt{done} events}, with per-segment \texttt{best\_val}
  from 3.6279 down to 3.5806 and back to 3.6109. The global minimum, \textbf{3.5806}, is \emph{better} than
  the 10k-step run's 3.621 that we declared the winner --- so that suite's ranking inverts
  depending on which number a reader takes, and we reported only the reading that supported our
  conclusion without knowing we were choosing. Token accounting for that run is also wrong: every
  \texttt{done} records 98,304,000 tokens for a run that reached step 19,990, where the correct figure is
  655,360,000. \textbf{We withdraw any horizon claim from that suite in either direction.}
\item
  \textbf{We reported a source as unlocatable while it sat in the repository.} \S7.3 item 2 previously
  stated that we could not find the artifact behind a ``29/29 oracle tests'' figure. It is at
  \path{experiment-notes/arch-metal/51-m5-128m-optimizer-funnel-preflight.md}, line 291.
\item
  \textbf{A ledger lost three completed jobs, and they were the ones that mattered.} The exact-128M
  LR spot-check ledger recorded one job as \texttt{"running"} and omitted three others entirely, although
  all four had finished on the same afternoon. Those three runs are what opened \S8.3.
\item
  \textbf{We broke our own reporting rule two sections after stating it.} \S3.2 says a \texttt{best\_val} field
  is not a paired snapshot and is never reported as a ranking. \S4.1 then reported four \texttt{best\_val}
  figures as end-of-run comparisons.
\item
  \textbf{Our knowledge base stopped recording results four weeks before we wrote this.} The project's
  \texttt{MASTER\_ARCHITECTURAL\_KB.md} changelog ended at 2026-07-20, so every suite in \S4 --- 120 runs at
  five seeds --- was absent from the document nominally serving as the project's verified record.
\item
  \textbf{A summary artifact reported one arm at 80\% of the horizon of the others.} In
  \texttt{crossover50m}, three of Mamba-2's five seeds ran on a 49.17M-token evaluation grid and two on
  a 49.99M one --- the only arm anywhere in these suites whose seeds end on different grids. The
  \texttt{summary.json} was generated by an aligner taking the \emph{exact intersection} of the two grids,
  which for that arm ends at \textbf{40.16M tokens}, and the figure it produced was 4.759; on the
  fixed aligner it is \textbf{4.656}. To be precise about scope, since we are auditing ourselves and
  not only the artifact: \textbf{this defect is not in the \S4.5 board.} That board's Mamba-2 row is
  \texttt{crossover50m\_matched32}, whose forty runs all end at 49,987,584 tokens, and it reproduces at
  4.596 on recomputation. The stale figure reached a \emph{different} rendering of this paper (see
  item 8). What the defect does touch here is \S4.5's comparison of arm movement against suite 22,
  which quotes suite-22 values measured at 49,594,368 tokens against suite-26 values at
  49,987,584 --- a 0.8\% horizon difference, now stated in that section.
\item
  \textbf{The source file of this manuscript was destroyed, and the reconstruction introduced a new
  error while fixing a real one.} A history rewrite of this repository stripped the Markdown
  source from the working tree and from every commit. The paper was rebuilt from the run
  artifacts. That rebuild correctly caught defect 7 above and a mis-stated parameter count, and
  correctly reported a normal-quantile interval bug in its own checking script. It also
  ``corrected'' the \S4.5 tie between attention and the 10$\times$minGRU + 2$\times$attention hybrid, on the
  grounds that the 4.232 figure ``appears in no surviving run''. That correction is itself wrong,
  and \S4.5 records why: 4.232 is the hybrid in \texttt{crossover50m\_matched32}, measured at the same
  49,987,584-token point as the attention row, while the replacement figure of 4.275 comes from
  \texttt{crossover50m}, where that arm's seeds end at 49,594,368 --- a different horizon. The
  reconstruction replaced a matched-horizon comparison with a mismatched one, in the same section
  where it documented horizon mismatch as the defect.
\item
  \textbf{A correction of ours was itself wrong, and a concurrent audit caught it.} We diagnosed the
  two champion BPB figures in \S7.2 as a digit transposition, because 2.015756 matched no artifact
  on disk. A separate audit of the same repository found \texttt{DECISIONS.md} M15 --- written before the
  audit pass we compared against --- independently recording 2.0158, which corroborates the figure
  as a third run's real value rather than a typo. Acting on our diagnosis would have deleted a
  real measurement from four documents. We had reasoned from absence of evidence, on exactly the
  axis this paper warns about, while writing the section that warns about it.
\end{enumerate}

The pattern across all nine is that \textbf{the failure mode is bookkeeping, not experimentation.} Runs
were executed carefully, gated, and seeded; they were then summarised into notes, ledgers, tables
and --- in the last case --- an entire replacement manuscript, and the summaries drifted from the
artifacts in ways invisible from the summaries alone. Recommendation 3 in \S9 (fingerprint every
job) addresses the input side of this. It does not address the output side, and on the evidence
above the output side is where our failures actually occurred. We now regard \emph{``re-derive every
reported number from the raw run directory before publication''} as a required step rather than a
diligence nicety, and \S10 states which numbers in this paper were and were not produced that way.

\begin{center}\rule{0.5\linewidth}{0.5pt}\end{center}

\section{8. The tuning confound: an alternative reading, and direct evidence that it bites}\label{the-tuning-confound-an-alternative-reading-and-direct-evidence-that-it-bites}

Every result in \S4 was produced under \textbf{standard parametrization} with a single global Muon
learning rate of 6e-4, held fixed across every mixer, every batch size and every schedule. This is
the largest uncontrolled factor in the paper, and it deserves its own section rather than a bullet
in \S7.2, because it admits an alternative reading of our central finding that we cannot currently
rule out.

The section runs in four steps: the competing explanation (\S8.1); what our seed structure does and
does not already bound (\S8.2); a measured instance of the same defect on our own optimizer axis,
where an inherited learning rate cost nearly twice the deciding margin (\S8.3); and the arm we
specify to settle it (\S8.4).

\subsection{8.1 The alternative explanation we cannot exclude}\label{the-alternative-explanation-we-cannot-exclude}

Our claim is that the crossing token moves with batch size and learning-rate horizon. A competing
explanation is narrower and less interesting:

\begin{quote}
A single global learning rate cannot be simultaneously near-optimal for a 12-layer attention
stack and a 12-layer minGRU stack, nor for the same stack at batch 8 and batch 32. What we
measured as ``the ranking depends on the recipe'' may be partly ``one arm was closer to its own
optimum than the other, and which arm that was changed with the recipe.''
\end{quote}

This is not a hypothetical concern. Three independent results make it concrete:

\begin{enumerate}
\def\labelenumi{\arabic{enumi}.}
\tightlist
\item
  \textbf{Optimal learning rate is horizon-dependent.} Bjorck et al.~show optimal LR falls as the token
  horizon grows and follows its own scaling law \cite{ref31}. Our suite 22/24 runs use a 50M-token cosine
  and our suite 23 run a 20M-token cosine; a single fixed peak LR of 6e-4 cannot be optimal for
  both, independently of the schedule-shape effect we attribute the 12.34M $\rightarrow$ 14.58M shift to.
\item
  \textbf{Optimal hyperparameters are batch-dependent in ways beyond the learning rate.} Marek et al.
  show that scaling to small batches correctly requires holding Adam's second-moment half-life
  fixed in tokens, not its decay rate \cite{ref30}. We changed batch size by 4$\times$ between suites 22 and 25
  and adjusted nothing.
\item
  \textbf{Parametrization determines how much of this matters.} Kalra and Barkeshli localize $\mu$P's
  advantage over standard parametrization under AdamW to maximizing the embedding-layer learning
  rate, which in standard parametrization acts as a bottleneck that induces instability \cite{ref35}. Our
  embedding layer is trained at the same global rate as everything else.
\end{enumerate}

If the competing explanation is right in full, the correct statement of our result weakens to
``under-tuned arms reorder unpredictably,'' which is a much less useful claim. If it is right only
in part --- our current expectation --- the crossing token still moves, but by less, and the honest
error bar on ``12.35M'' widens beyond the seed-derived interval we report.

\subsection{8.2 What the existing evidence does and does not already constrain}\label{what-the-existing-evidence-does-and-does-not-already-constrain}

Two things partially bound the concern, and we state them without overclaiming.

\textbf{Sign consistency across seeds does not help here.} All five seeds agreeing on the sign of the
gap constrains \emph{seed} noise. It says nothing about a systematic mis-tuning that is shared by all
five seeds of an arm, which is exactly what a wrong global learning rate would be.

\textbf{The 6e-4 value is not arbitrary, but it was tuned on neither arm at these recipes.} It was
inherited from the earlier 3070 Ti suites. That makes it, at best, a value tuned for one
architecture at one batch size at one horizon --- precisely the transfer that \cite{ref31} and \cite{ref30} say
should not be assumed.

\subsection{8.3 The concern is not hypothetical: the ordering it hides crosses over}\label{the-concern-is-not-hypothetical-the-ordering-it-hides-crosses-over}

The argument in \S8.1 was an appeal to the literature. We can do better. Our own optimizer axis
contains the failure \S8.1 describes, and measuring it properly produced a result stronger and
stranger than the concern we set out to test.

The \S5 funnel selected \texttt{muon\_polar\_adamw} at a matrix learning rate of \textbf{0.05} over
\texttt{normuon\_adamw} at \textbf{0.1}, by 0.031226 BPB. Both values were tuned at \texttt{arch02-16m} and \textbf{neither
was re-tuned at 128M}, the scale at which the selection was made. We re-ran both candidates across
a matched eight-point learning-rate grid at the \emph{exact} protocol that produced the selection ---
\texttt{arch02-128m}, 1,000 steps, seeds 42, 2026 and 1337, argv byte-identical to the recorded
\texttt{exact\_128m\_1000} jobs apart from the output path. \textbf{Fifty-two jobs, about 30 GPU-hours.}

{%
{\small\begin{longtable}[]{@{}
  >{\raggedright\arraybackslash}p{(\linewidth - 8\tabcolsep) * \real{0.2000}}
  >{\raggedright\arraybackslash}p{(\linewidth - 8\tabcolsep) * \real{0.2000}}
  >{\raggedright\arraybackslash}p{(\linewidth - 8\tabcolsep) * \real{0.2000}}
  >{\raggedright\arraybackslash}p{(\linewidth - 8\tabcolsep) * \real{0.2000}}
  >{\raggedright\arraybackslash}p{(\linewidth - 8\tabcolsep) * \real{0.2000}}@{}}
\toprule\noalign{}
\begin{minipage}[b]{\linewidth}\raggedright
matrix LR
\end{minipage} & \begin{minipage}[b]{\linewidth}\raggedright
\texttt{muon\_polar\_adamw}
\end{minipage} & \begin{minipage}[b]{\linewidth}\raggedright
\texttt{normuon\_adamw}
\end{minipage} & \begin{minipage}[b]{\linewidth}\raggedright
gap (Polar - NorMuon)
\end{minipage} & \begin{minipage}[b]{\linewidth}\raggedright
leader
\end{minipage} \\
\midrule\noalign{}
\endhead
\bottomrule\noalign{}
\endlastfoot
0.0035 & 2.138058 & 2.144220 & -0.006162 & \textbf{Polar} \\
0.005 & 2.125631 & 2.128408 & -0.002778 & \textbf{Polar} \\
0.008 & \textbf{2.125430} & \textbf{2.121334} & +0.004097 & NorMuon \\
0.0125 & 2.139372 & 2.121866 & +0.017505 & NorMuon \\
0.018 & 2.141161 & 2.129755 & +0.011406 & NorMuon \\
0.025 & 2.152357 & 2.135230 & +0.017127 & NorMuon \\
0.035 & 2.158727 & 2.146234 & +0.012493 & NorMuon \\
0.05 & 2.171450 & 2.155408 & +0.016042 & NorMuon \\
\end{longtable}}
}

Means over three seeds; gaps paired by seed. \textbf{Every row is sign-consistent across all three
seeds} --- twenty-four paired comparisons, none dissenting from its side.

\textbf{The ordering crosses over between lr 0.005 and lr 0.008.} \texttt{muon\_polar\_adamw} leads at the two
lowest learning rates, \texttt{normuon\_adamw} at the six highest, each unanimously. The mechanism is
visible in the table: the two have \textbf{offset flat basins} --- Polar's minimum region spans
0.005--0.008, NorMuon's spans 0.008--0.0125 --- so each wins inside its own and loses inside the
other's. Neither basin has a resolvable interior point at n = 3; both are flat to within the
per-cell seed spread of roughly 0.006.

\textbf{There is therefore no recipe-independent answer to which of these two optimizers is better at
this scale.} The question the funnel asked does not have one.

What the funnel did was pick a side of that crossing without knowing it existed. It compared Polar
at 0.05 against NorMuon at 0.1 --- both far up the high-LR wall, in the region where NorMuon leads.
Measured against their own optima at this protocol, the inherited values cost:

\begin{verbatim}
muon_polar_adamw   selected 0.05, best 0.008   penalty 0.046020 +/- 0.003175   = 1.47x the margin
normuon_adamw      selected 0.1,  best 0.008   penalty 0.080400 (n = 2)        = 2.57x
\end{verbatim}

The Polar figure is a paired three-seed interval that excludes zero. Both inherited learning rates
were 6--12$\times$ too high, unequally so, and the difference between those two penalties exceeds the
margin that decided the outcome.

\textbf{What we claim.} The selection recorded in \S5 is \textbf{retired}: \path{research/champion-run.json}
remains \texttt{"locked":\ false}. We do \textbf{not} crown \texttt{normuon\_adamw}; that would repeat the original
error with the sign flipped, since it loses at the low end and at each candidate's own best cell
the gap is +0.004097 $\pm$ 0.005085, which spans zero.

\textbf{Limits.} Sign tests carry this result: only three of eight matched cells have an interval
excluding zero, because the effects sit at or below the per-cell seed spread. Neither optimum is a
located point. This is one protocol --- 128M, 1,000 steps, Metal on an M5 Pro. Four cells were
re-run against their recorded funnel values as a drift check and moved by -0.0019, +0.0003,
-0.0058 and +0.0021, at most 0.26\%, far below the effects above.

\textbf{What this section is really reporting.} \S4 shows a ranking between two architectures crossing
over as a function of \emph{token budget}. This section shows a ranking between two optimizers crossing
over as a function of \emph{learning rate} --- the same structure, on an axis we had treated as a nuisance
parameter rather than a variable. The funnel's error was not that it tuned badly. It was that it
reported a ranking measured at one point on an axis along which the ranking reverses.

The grid history is part of the result and we state it rather than presenting the final design as
if it were the plan. This began as twelve jobs to re-tune one learning rate. It became fifty-two,
across five rounds, because each round revealed the previous one had compared something unequal: a
truncated grid, then a boundary minimum, then one candidate given lower learning rates than the
other. Three times we drew a conclusion from two seeds that a third seed overturned. That
trajectory is the honest cost of controlling this confound once, on one axis, for two candidates,
and it is the strongest argument we can make that the confound is not second-order.

Full grid, per-seed values, drift checks and limits: \path{research/d7-lr-retune.json}, regenerated by
\path{scripts/d7_analyze.py} from the run ledger rather than transcribed. Runner:
\path{scripts/d7_lr_retune.py}.

\subsection{8.4 The follow-up, specified}\label{the-follow-up-specified}

We name the arm rather than gesture at it, so that it can be run or criticized.

\textbf{Design.} Repeat the suite 24 recipe --- 12L/768d, batch 32, \texttt{eval\_iters\ =\ 20}, 50M-token cosine
horizon, 20M-token stop, seeds \texttt{1337,\ 42,\ 100,\ 2026,\ 777} --- as a 2 $\times$ 2:

{%
{\small\begin{longtable}[]{@{}
  >{\raggedright\arraybackslash}p{(\linewidth - 4\tabcolsep) * \real{0.3333}}
  >{\raggedright\arraybackslash}p{(\linewidth - 4\tabcolsep) * \real{0.3333}}
  >{\raggedright\arraybackslash}p{(\linewidth - 4\tabcolsep) * \real{0.3333}}@{}}
\toprule\noalign{}
\begin{minipage}[b]{\linewidth}\raggedright
\end{minipage} & \begin{minipage}[b]{\linewidth}\raggedright
attention
\end{minipage} & \begin{minipage}[b]{\linewidth}\raggedright
minGRU
\end{minipage} \\
\midrule\noalign{}
\endhead
\bottomrule\noalign{}
\endlastfoot
\textbf{standard parametrization, global LR} & suite 24 (already run) & suite 24 (already run) \\
\textbf{$\mu$P, per-arm LR swept at a narrow proxy width} & new & new \\
\end{longtable}}
}

with the $\mu$P cells parametrized per \cite{ref12} and the proxy sweep run at a reduced width, transferring
the tuned LR to the 768-dim target. Depth is held fixed at 12 layers throughout, which sidesteps
the depthwise limitations Depth-$\mu$P identifies for stacks whose residual blocks are themselves deep
\cite{ref33}. Two secondary arms are worth the marginal cost:

\begin{itemize}
\tightlist
\item
  \textbf{A per-layer standard-parametrization prescription} as in Everett et al.~\cite{ref32}, since their
  result is that this can outperform $\mu$P and that all parameterizations admit transfer --- if the
  crossing token is stable under \emph{both} $\mu$P and per-layer SP but not under a global LR, the finding
  is about tuning quality, not about $\mu$P specifically.
\item
  \textbf{An embedding-layer-LR-only ablation} as in Kalra and Barkeshli \cite{ref35}: raise only the embedding
  learning rate by a factor of width under standard parametrization. This is the cheapest possible
  probe, and if it moves the crossing token materially it localizes the confound to one layer.
\end{itemize}

\textbf{Pre-registered readouts.} The follow-up answers a yes/no question, and we state the decision
rule before running it:

{%
{\small\begin{longtable}[]{@{}
  >{\raggedright\arraybackslash}p{(\linewidth - 2\tabcolsep) * \real{0.5000}}
  >{\raggedright\arraybackslash}p{(\linewidth - 2\tabcolsep) * \real{0.5000}}@{}}
\toprule\noalign{}
\begin{minipage}[b]{\linewidth}\raggedright
Outcome
\end{minipage} & \begin{minipage}[b]{\linewidth}\raggedright
Reading
\end{minipage} \\
\midrule\noalign{}
\endhead
\bottomrule\noalign{}
\endlastfoot
Both $\mu$P crossings land within the suite 22/24 per-seed bands (early 1.03--1.09M, late 11.93--12.58M) & Recipe dependence is not a tuning artifact. \S4.3's schedule effect stands as stated. \\
Crossings move but the \textbf{ordering of recipes} is preserved --- 20M cosine still crosses later than 50M cosine & The schedule effect is real; its magnitude in \S4.3 is inflated by mis-tuning and should be restated as an upper bound. \\
$\mu$P crossings collapse onto a single token across schedules & The schedule effect is a tuning artifact. \S4.3 must be withdrawn and \S4.4's batch result re-examined. \\
The batch-8 arm (suite 25) develops a crossing under $\mu$P within 7.4M tokens & The batch effect is a tuning artifact, and suite 14's original 6.6--7.4M window is partially rehabilitated. \\
\end{longtable}}
}

\textbf{Cost.} Twenty jobs at 20M tokens for the main 2 $\times$ 2, plus a proxy-width LR sweep, plus ten jobs
for each secondary ablation. On the GH200 configuration used for suites 23--25 this is roughly the
cost of suite 24 plus suite 25 together --- a small fraction of the 50-job suite 22 grid, and the
highest-value single experiment remaining in this line of work.

\textbf{Until it is run}, every crossing token in this paper should be read as ``the crossing token for
this architecture pair \emph{at a global 6e-4 Muon learning rate under standard parametrization},'' and
the \S9 recommendations should be read as conditional on the same. We consider recommendation 2
(pin the schedule when truncating) robust to the outcome, because it is a statement about
comparing a prefix to its parent recipe rather than about either being correctly tuned; we
consider the specific figure ``2.2M tokens, 18\% of the location'' provisional.

\begin{center}\rule{0.5\linewidth}{0.5pt}\end{center}

\section{9. Recommendations}\label{recommendations}

These are the rules we now apply, each earned by a specific failure above.

\begin{enumerate}
\def\labelenumi{\arabic{enumi}.}
\tightlist
\item
  \textbf{Report the crossing, not just the endpoint.} If two curves cross, the crossing token is a
  result and belongs in the paper with per-seed values. If they cross twice, say so.
\item
  \textbf{Never truncate a run without pinning the schedule.} Fix the learning-rate horizon to the
  \emph{reference} budget and stop early, rather than letting the cosine follow the truncated step
  count. Otherwise the prefix is a different recipe (\S4.3). Record \texttt{lr\_max\_steps} alongside
  \texttt{max\_steps} in every config, and verify it per job before reading results.
\item
  \textbf{Fingerprint the recipe on every job and verify before aggregating.} Suite 22's mixed
  batch / \texttt{eval\_iters} board would have been read as a bakeoff without that check. At minimum:
  batch size, evaluation count, max steps, LR horizon, compile flag, and kernel version.
\item
  \textbf{Never table cells from different recipes together.} Even when the token markers match.
\item
  \textbf{Five seeds for a sign, two for nothing.} A two-seed 95\% interval carries a t$_1$ = 12.706
  multiplier and will almost never separate close candidates. If you only have two seeds, report
  both values and claim only a consistent ordering.
\item
  \textbf{A short learning-rate screen is not a learning-rate selection.} Screen at a horizon within
  the same order of magnitude as the target, or accept that you are ranking EMA lag.
\item
  \textbf{A minimum-over-evaluations field is not a ranking.} Compare paired snapshots at matched
  token counts.
\item
  \textbf{Record systems-blocked arms as exclusions with the blocking reason}, never as silent
  substitutions or omissions.
\item
  \textbf{Assume no safe direction of extrapolation.} Small-scale screens under-rank the eventual
  large-scale champion \emph{and} the large-scale champion under-performs at small scale (\S5.5). Both
  directions failed here.
\end{enumerate}

\begin{center}\rule{0.5\linewidth}{0.5pt}\end{center}

\section{10. Reproduction and artifacts}\label{reproduction-and-artifacts}

\textbf{Repository.} \url{https://github.com/bharathvbcr/MachineLearning} (public). Every path below is
relative to that repository root, at commit \texttt{a83257a} or later. Two directories referenced by
these instructions are deliberately not committed because they are build outputs with in-repo
generators: \path{nanolab/out/} (raw per-job run directories) and
\path{Rust_MLKit/arch_02_value_resid/metal-native/golden/fwd/} (parity activations; regenerate with
\path{python3\ metal-native/scripts/export_goldens.py\ -\/-out\ metal-native/golden}).

All mixer suites run through one entry point:

\begin{verbatim}
python -m nanolab.crossover_replicate list
python -m nanolab.crossover_replicate smoke            # 40-step GPU/CPU check
python -m nanolab.crossover_replicate launch --workers 4
python -m nanolab.crossover_replicate isolates         # matched20 -> bs8 -> matched32
python -m nanolab.crossover_replicate status --out <run-dir>
python -m nanolab.crossover_replicate timing --out <run-dir> --json
python -m nanolab.crossover_replicate plot
\end{verbatim}

\textbf{A note on timing.} No suite in this paper reports GPU hours, because the trainer computed the
elapsed time of every run, printed it, and discarded it --- the \texttt{done} record carried only
\texttt{best\_val} and \texttt{tokens}. That is fixed as of 2026-08-22 (\texttt{Logger.done} now persists \texttt{elapsed\_s}),
but it is not retroactive: the \texttt{timing} command back-estimates suites 22--26 from median \path{tok/s}
and labels those rows \texttt{estimated\_from\_median\_tok\_s}, separately from \texttt{measured}. Treat any GPU-hour
figure for suites 22--26 as an estimate.

Per-suite output roots, each containing \texttt{metrics.jsonl} + \texttt{config.json} per job plus \texttt{queue.json}:

{%
{\small\begin{longtable}[]{@{}llll@{}}
\toprule\noalign{}
Suite & Output root & Job prefix & Jobs \\
\midrule\noalign{}
\endhead
\bottomrule\noalign{}
\endlastfoot
22 & \path{nanolab/out/crossover50m/} & \texttt{cx50\_} & 50 \\
23 & \path{nanolab/out/crossover20m_locked/} & \texttt{cx20\_} & 10 \\
24 & \path{nanolab/out/crossover20m_matched_lr/} & \texttt{cx20h\_} & 10 \\
25 & \path{nanolab/out/crossover8m_bs8/} & \texttt{cx8\_} & 10 \\
26 & \path{nanolab/out/crossover50m_matched32/} & \texttt{cx32\_} & 40 \\
\end{longtable}}
}

Locked summary statistics --- the numbers in \S4 are read from these, not retyped from logs:

\begin{itemize}
\tightlist
\item
  \path{experiment-notes/nanolab/artifacts/22-suite22_lock.json}
\item
  \path{experiment-notes/nanolab/artifacts/23-locked20_lock.json}
\item
  \path{experiment-notes/nanolab/artifacts/24-matched20_lock.json}
\item
  \path{experiment-notes/nanolab/artifacts/25-bs8_lock.json}
\item
  \path{experiment-notes/nanolab/artifacts/26-matched32_lock.json}
\end{itemize}

Figures: \texttt{experiment-notes/nanolab/figures/\{22-attn-mingru-crossover,\ 23-locked20-crossover,\ 24-matched20-crossover,\ 25-bs8-crossover,\ 26-matched32-ranking\}.\{png,svg,pdf\}}.

Optimizer funnel: manifest \path{research/optimizer-study.json} (candidate list, stage definitions,
hard gates, required metrics, and the two blocked-candidate reasons); results
\path{research/native-optimizer-funnel.json} (106 jobs); champion record \path{research/champion-run.json};
exact-resume gates \path{research/exact-128m-gate-polar.json} and \path{research/exact-128m-gate.json}.

Full lab notes with per-suite failure logs, evidence grades and interpretation boundaries:
\path{experiment-notes/nanolab/14}, \texttt{15}, \texttt{16}, \texttt{22}, \texttt{23}, \texttt{24}, \texttt{25}, \texttt{26}.

\textbf{Non-nanolab sources this paper depends on.} An earlier version of this section listed only the
\path{nanolab/} and \path{research/} paths above, while \S5 and \S7.3 in fact rest on the Metal/Rust tracks.
Those paths are:

{%
{\small\begin{longtable}[]{@{}
  >{\raggedright\arraybackslash}p{(\linewidth - 2\tabcolsep) * \real{0.5000}}
  >{\raggedright\arraybackslash}p{(\linewidth - 2\tabcolsep) * \real{0.5000}}@{}}
\toprule\noalign{}
\begin{minipage}[b]{\linewidth}\raggedright
Used by
\end{minipage} & \begin{minipage}[b]{\linewidth}\raggedright
Path
\end{minipage} \\
\midrule\noalign{}
\endhead
\bottomrule\noalign{}
\endlastfoot
\S5.1 funnel design, stage definitions, the 29/29 orchestration oracle & \path{experiment-notes/arch-metal/51-m5-128m-optimizer-funnel-preflight.md} \\
\S5.5 downward-transfer table (Polar vs NS5 at \texttt{sota}, 3k, seeds 1337/42) & \path{experiment-notes/arch-metal/53-loop-research-soft-fatiled-2026-07-23.md} and its machine-readable board \path{Rust_MLKit/arch_02_value_resid/metal-native/out/loop_research_2026-07-23/RESULTS.json} \\
\S7.3 withdrawal 2 (the retracted Metal-vs-CUDA BPB comparison) & \path{experiment-notes/arch-metal/50-*}, \path{Rust_MLKit/arch_02_value_resid/metal-native/README.md}, \path{metal-native/DECISIONS.md} \\
\S7.3 withdrawal 3 (the retracted \textasciitilde238k/\textasciitilde240k tok/s figures and their re-baseline) & withdrawn values in \path{experiment-notes/arch-metal/52-context-crossover-metal.md}; honest re-baseline in \path{\ldots/53-\ldots}, \S``mixer honesty'' \\
\S7.3 withdrawal 5 (62--100\% of \textasciitilde273 GB/s; \textasciitilde77\% dispatch overhead) & \path{Rust_MLKit/gemma-metal/bench/results/kernel_roofline_finding.json}, \path{Rust_MLKit/gemma-metal/docs/bottleneck.md} \\
\S8 LR-transfer evidence (new; see \S8.2) & \path{research/champion-run.json} $\rightarrow$ \texttt{lr\_transfer\_finding}; raw runs in \path{out/funnel/polar_exact_lr_spot/} (not committed --- \path{out/} is a build output; \texttt{ledger.json} in that directory carries the reconciled record) \\
\end{longtable}}
}

\textbf{Two scope caveats on the re-baselined throughput figures.} The 144,367 / 66,382 / 33,321 tok/s
values in \S7.3 are Metal measurements at sequence length 256 on 0.78M--1.05M-parameter models, over
40 steps. Every quality result in \S4 is CUDA at sequence length 512 on 124M-parameter models. ``The
same three arms'' refers to the mixer \emph{type} only; the two sets of numbers share no shape, scale,
backend or sequence length and must not be placed in one table.

\textbf{What the lock files do and do not contain.} \path{experiment-notes/nanolab/artifacts/*_lock.json}
carry the crossing points, the 19.677M gap, and the 50M ranking board (\texttt{arm}, \texttt{mean}, \texttt{lo}, \texttt{hi}).
They do \textbf{not} carry the per-marker tables of \S4.2--\S4.4, their confidence intervals, or the
per-seed ranges shown in the \S4.5 board. Those are recomputed from the per-job \texttt{metrics.jsonl}
under \path{nanolab/out/}, which is not committed (see the note on build outputs above, and gap B3 in
\path{docs/ISSUES_AND_GAPS_2026-08-22.md}). Reproducing the \S4 tables therefore requires re-running the
suites or obtaining the run directories; reproducing the crossings and the 50M board does not.

Two reproduction caveats, stated rather than buried:

\begin{itemize}
\tightlist
\item
  \textbf{Suites 14, 15 and 16 have no preserved replay command.} The closest runner is
  \texttt{python\ -m\ nanolab.train}; every recorded field in the relevant \texttt{config.json} must be
  reconstructed before comparing outputs. This is why those suites are cited as observations and
  never as the basis for a shipped claim.
\item
  \textbf{Suite 23's artifacts are the short-cosine run and were not regenerated.} The runner now sets
  \texttt{CROSSOVER\_LR\_HORIZON=50000000} so a future launch preserves the long cosine; reproducing
  suite 23 as recorded requires the old behavior, and a schedule-matched prefix must use a new
  output directory.
\end{itemize}

\begin{center}\rule{0.5\linewidth}{0.5pt}\end{center}

\section{11. Conclusion}\label{conclusion}

We tried to replicate a crossover we had measured ourselves, at five seeds instead of one, on
better hardware, and over six times the token budget. The shape survived. The number did not.

The token at which attention overtakes minGRU on this model is 12.35M under a 50M-token cosine,
14.58M under a 20M-token cosine, and does not exist within 7.4M tokens at batch size 8 --- and the
originally reported 6.6--7.4M window turns out to sit inside a region where minGRU leads by
0.13--0.18 nats on every seed of every batch-32 recipe we ran. On an entirely separate axis, an
optimizer that won a two-seed 16M-parameter stage finished last of four at 128M, and the eventual
128M champion loses when run back down to small scale.

The uncomfortable version of the lesson is that the cheap screen cannot be trusted even as a
shortlist. The useful version is narrower and actionable: \textbf{the ranking is a property of the
measurement recipe as much as of the method, so the recipe must be pinned, fingerprinted,
reported, and never silently varied between the screen and the run it is screening for.}

\begin{center}\rule{0.5\linewidth}{0.5pt}\end{center}

\subsection{Provenance of this draft}\label{provenance-of-this-draft}

All 38 literature identifiers cited in \S1, \S2 and \S8 were resolved through the ScholarLM / WisDev
MCP (\texttt{wisdevSearchPapers}, \texttt{wisdevPaperLookup}) against arXiv, OpenAlex, Semantic Scholar and
Crossref. Every arXiv ID in the reference list was fetched individually and its title, authors,
year and abstract confirmed against the claim made about it in the text; the discovery searches
that surfaced \cite{ref29}--\cite{ref32} and \cite{ref35} were run in \texttt{methodology} research mode. No citation was written
from recall, and no claim is attributed to a paper whose abstract was not read.

Every experimental number in \S4 and \S5 is read from the locked artifacts listed in \S10. Suite
statistics come from the \texttt{*\_lock.json} files; optimizer-funnel figures are recomputed directly
from \path{research/native-optimizer-funnel.json}, which is also the source of the interval correction
in \S5.3.

ScholarDoc's drafting pipeline was \textbf{not} used for the prose. Its retrieval tools were, as
described above. The local stdio server initially returned only a section scaffold because the
MCP child process was launched without \texttt{GOOGLE\_CLOUD\_PROJECT} and so could not reach the Vertex
backend that the orchestrator was already using; that was fixed, and DocGen was then run three
times against a healthy backend. All three runs failed to ground: the best grounded on a single
source on an unrelated topic (information-retrieval ranking-model distillation, reached by
semantic drift on the word ``ranking''), and the last grounded on \textbf{zero} sources while still
emitting a seven-section draft and a numeric review score. None of that output was merged. The
tooling defects are recorded in \path{docs/ISSUES_AND_GAPS_2026-08-22.md} \S1.

